\begin{document}
 	
Powered by Panopto
2025-26 - C004205A_2025 - C004205A - Chemie
Hoorcollege

Auto-generated captions may contain errors.

Werken? Ja, ik stel voor dat we van start gaan.
1:03
Dus we hebben die thermodynamica opgebouwd.
1:10
Samenvatting van evenwicht principe is dat de vrije enthalpie van een systeem onder PT controle zo laag mogelijk wordt.
1:16
We hebben een uitdrukking voor DG in termen van een CDT, een VDP en een sommatie over alle componenten in het systeem.
1:25
Nu Sidney. Waarbij onder PT controle eigenlijk enkel die laatste bijdrage overblijft waarmee het systeem zijn vrije enthalpie kan minimaliseren.
1:34
We hebben dat toegepast op. Een twee fasen evenwicht waarbij we een zuivere stof hebben die als vloeistof of als gas kan voorkomen.
1:46
Bijvoorbeeld wat we hebben vastgesteld is dat die relatie gaat aanleiding geven tot of kan vertaald worden in het feit dat de chemische
1:56
potentiaal van die component A in de vloeistof fase moet gelijk zijn aan de chemische energie van die component A in de gasfase.
2:04
Dat is fysisch evenwicht. Als ik een component kan uitwisselen tussen beide fasen,
2:11
gaat evenwicht ontstaan als de chemische potentiaal van die component in beide fasen dezelfde is.
2:16
Op basis daarvan kan je bijvoorbeeld de dampdruk gaan uitwerken, de dampdruk van zo'n idee die je hier hebt.
2:21
Als je evenwicht hebt tussen de damp en de vloeistof fase kan je gaan uitdrukken in termen van gestabiliseerde fysieke chemische grootheden.
2:29
Ik heb daarvoor een aantal solventen of oplosmiddelen data verzameld.
2:37
Hier is de stok. Ja, en ik kan er eentje uitpikken, bijvoorbeeld koolstof nitroglycerine.
2:44
Dat is één van die componenten waarvan je de naamgeving toch moet beheersen.
2:51
Die zijn koolstof gebonden aan vier. Chloor is een beetje analoog aan methaan,
2:55
maar elke waterstof is vervangen door een chloor en krijg je data rond vorming enthalpie van CL vier in de vloeistof fase vorming
3:00
van special vier in de gasfase en je krijgt ook waarde voor de absolute entropie voor special vier als vloeistof en als gas.
3:09
Hoe halen we daar de dampdruk van Special vier uit bij een bepaalde temperatuur?
3:18
Dat zijn data bij twee 98 Kelvin. Stel dat ik wil weten wat de dampdruk van K4 is bij 320 Kelvin.
3:23
Hoe kan ik met die gegevens aan de slag? Om die dampdruk te bepalen.
3:31
En wat waren de verschillende dingen waar we rekening mee moesten houden?
3:53
Als je die relatie uitwerkt. Wat kan je doen met die chemische middelen van de zuivere vloeistof?
4:01
Het was die ene slide die je. Uit je hoofd moest leren.
4:17
Blijkbaar is dat een proces dat nog. Bezig is.
4:23
Drie vakjes. Moet haalbaar zijn.
4:29
Voor een zuivere vloeistof kunnen we de druk afhankelijkheid laten vallen.
4:34
Omdat het molaire volume zo klein is, kunnen we altijd gelijkstellen aan de standaardwaarde.
4:37
Standaardwaarde betekent De druk is één bar.
4:42
De temperatuur is nog altijd een variabele voor de gasfase als het een ideaal gas is, dat was het middelste vakje van die slide,
4:45
kan je de chemische potentiaal schrijven als een combinatie van de tabel en
4:53
de standaardwaarde en dan druk afhankelijkheid van de vorm RTL partieel druk
4:57
gedeeld door totaal druk kan je afleiden uit de ideale gaswet in combinatie met die uitdrukking voor DG en dan gefocust op de druk afhankelijkheid.
5:03
Dat betekent dat je de evenwichts druk kan schrijven als standaard druk maal de exponentiële van min.
5:13
Als ik dan maar het andere lid breng. Gemis van de standaardwaarde A in de gasfase.
5:23
Standaardwaarde van A in de vloeistof fase gedeeld door een bedrag R2.
5:29
Wat we kunnen samenvatten dat verwijst allebei naar de zuivere stof A Als gas en als vloeistof
5:35
is dus de verdamping vrije enthalpie min delta verdamping vab g nul gedeeld door R2 En pas op,
5:42
dit is een grootheid die afhangt van de temperatuur.
5:52
De nul wijst enkel naar het feit dat het een waarde is bij een druk van één bar.
5:56
Dus als we de dampdruk van zo'n oplosmiddel willen bepalen in eenheden van standaard druk, dan moeten we weten wat de verdamping is.
6:01
Vrije enthalpie De standaard verdamping is vrije enthalpie.
6:10
Bij de gegeven temperatuur is van die vloeistof.
6:13
Als de temperatuur 23 Kelvin is en de data zijn gegeven twee 98 Kelvin.
6:19
Hoe konden we aan de slag gaan om die delta geen nul te bepalen? Op basis van de definitie delta geeft nul.
6:25
Bij de gegeven temperatuur is delta H02 die temperatuur -2 delta S0 bij die temperatuur.
6:36
Er zijn nog altijd dingen die we niet uit de tabel rechtstreeks kunnen aflezen,
6:46
want dit gaat ook over de delta H bij 23 Kelvin en dat gaat ook over de delta S bij 320 Kelvin.
6:50
Hoe kunnen ze die in verband brengen met de data die in die tabel staan?
6:57
Een delta is het verschil tussen de enthalpie van de gasfase en de vloeistof fase.
7:12
Dat zijn dus verschillen tussen het getal dat je hier hebt en het getal dat je daar hebt.
7:19
Het zal dus 31 kilo per mol zijn voor chloroform voor CTL vier.
7:24
Als we daarmee werken is dat die twee 72,8 -317.
7:29
Rap uitgerekend is dat 44,2 of vijf 44,2 kilo per mol.
7:34
De Delta s. Dat is het verschil tussen waarden die in de gasfase gestabiliseerd zijn en die in de vloeistof fase getallenleer zijn.
7:47
Dat gaat 125,3 joule per Kelvin per mol zijn.
7:54
Pas op als we dingen berekenen. Dat hebben we vorige les gezien.
8:03
Je vult het dier niet in in kilo's. Je vult dat in in Zoelen per mol.
8:06
Dus dat gaat over 44.200 joule per mol gaan.
8:10
Maar dit is de waarde bij 320 Kelvin was een waarde met twee 98 Kelvin.
8:14
Wat konden we doen of wat hebben we gezegd over die temperatuurs afhankelijkheid?
8:19
Van Delta en Delta is. Eindelijk.
8:25
Die kunnen we als temperatuurs onafhankelijk beschouwen. Er wordt ingeprikt op mijn verhaal.
8:31
Dat is heel goed en dat is omdat we het verschil in warmtecapaciteit tussen wat je hier producten en reagentia gasfase en vloeistof fase noemt,
8:36
dat dat doorgaans relatief klein is.
8:44
In die zin een temperatuurs afhankelijkheid van zo'n delta wordt bepaald door het verschil in warmtecapaciteit tussen producten en reagentia.
8:46
Analoog een temperatuurs afhankelijkheid van een delta s wordt bepaald door het verschil in warmtecapaciteit gedeeld door de temperatuur.
8:57
Als we werken over als we kijken naar temperatuurs intervallen van een aantal tientallen Kelvin en die delta C.
9:06
Dat verschil in warmtecapaciteit is relatief klein en over dat soort beperkte temperatuurs
9:12
intervallen kunnen we de delta en delta als temperatuurs onafhankelijk beschouwen.
9:17
Nooit een delta G als temperatuurs onafhankelijk beschouwen.
9:22
Je hebt hier altijd die vermenigvuldiging twee delta s zitten waar je bij benadering kan doen is de delta en de delta is delta H0.
9:25
Dat is nul als temperatuurs onafhankelijk beschouwen. Dat doe je als je niet extra als je geen extra data ter beschikking hebt.
9:33
En de fout die je daar maakt is is doorgaans niet al te niet al te belangrijk.
9:40
En dus door die waarden in te vullen, te combineren met de temperatuur bepaal je de delta geen nul bij die specifieke temperatuur.
9:45
En dan kan je met deze relatie de dampdruk gaan uitrekenen en dan kan je eventueel
9:53
doen bij verschillende temperaturen en dan krijg je die cosinus klap klapper. Een relatie LNB versus één op drie die we voor CO2 getoond heb,
9:58
maar die je evengoed voor al die verschillende oplosmiddelen kan gaan en kan gaan bepalen.
10:04
Vandaag wil ik het niet verder hebben over dit soort vage evenwichten.
10:10
Ik ga ervan uit dat je dat beheerst. Het is daarom dat we er nog kort even zijn doorgegaan.
10:13
Ik wil de stap zetten naar chemische evenwichten.
10:20
Wat betekent dat we niet alleen stoffen kunnen hebben die tussen twee fasen worden uitgewisseld,
10:25
maar dat we ook praten over chemische reacties, evenwicht bij chemische reacties waarbij bindingen gevormd en gebroken kunnen worden.
10:31
Voordat we naar die chemische reacties gaan kijken, wil ik kort even terugkeren op de manier waarop we dat evenwicht hier hebben ingesteld.
10:46
En we komen aan die gelijkheid van chemische potentialen door te kijken hoe uitwisseling van A tussen beide fasen de vrije enthalpie verandert.
10:54
En dan kom je bij een vrije enthalpie verandering die geschreven kan worden als chemisch potentiaal van A in de gasfase
11:08
min chemische potentiaal van A in de vloeistof fase maal het aantal mol dat uitgewisseld wordt tussen beide fasen.
11:14
Als ik hier een aantal modi en weg neem, vermindert de vrije enthalpie met een bedrag chemische potentiaal van A in de vloeistof fase maal DN
11:22
en als ik dat daar terugplaatst dan stijgt de vrije enthalpie met een bedrag van de schaal van één.
11:32
De gasfase maal DN. Essentieel om een minimum van de vrije enthalpie te bereiken is dus dat er uitwisseling is van.
11:38
Stof van materie tussen beide fasen.
11:50
Als je die uitwisseling onmogelijk maakt, dan heeft ook die thermodynamisch analyse van evenwicht geen betekenis.
11:54
Het systeem verlaagt de vrije enthalpie door materie uit te wisselen tussen beide fasen.
12:00
We noemen dat een dynamisch evenwicht. Dat betekent ook als we in evenwicht zijn.
12:08
Als de chemische potentialen gelijk zijn, moet je niet denken dat er helemaal niets meer gebeurt.
12:13
Er gaan ook altijd moleculen verdampen van de vloeistof fase naar de gasfase.
12:18
Daar gaan ook altijd moleculen condenseren vanuit de gasfase in de vloeistof fase.
12:23
Maar die twee flux, die zijn even groot, die houden elkaar in evenwicht.
12:28
En aan die evenwichtssituatie wijzigt er verder niets. Dat noemen we een dynamisch evenwicht.
12:32
Het is een stationaire toestand, tijd onafhankelijk.
12:38
Maar onderliggend heb je eigenlijk twee stromen.
12:41
Er is verdamping en er is condenseren en die houden elkaar in balans.
12:44
Zonder zo'n dynamisch evenwicht heeft het geen zin om over thermodynamica te praten.
12:50
Denk aan de manier waarop die vrije enthalpie geminimaliseerd wordt.
12:55
Dat is door juist materiaal uit te wisselen tussen beide fasen.
13:00
Als die uitwisseling niet mogelijk is, dan kan dat systeem ook zijn vrije enthalpie niet gaan minimaliseren.
13:05
Dus dat idee van dynamisch evenwicht is eigenlijk essentieel in een thermodynamische analyse.
13:11
En dan zie je ook verschijnen als je naar chemische reacties gaat kijken.
13:16
Hier is een voorbeeld is een beetje een hypothetische reactie die getekend wordt,
13:20
maar je kan denken aan een reactie waarbij twee moleculen hier geschreven als A2 en B2,
13:26
dat zo stikstof en zuurstof die stikstof en zuurstof kunnen zijn en die kunnen reageren ter vorming van a b.
13:32
Dat zou dan stikstofmonoxide kunnen zijn en je kan ofwel wat er in links staat als de reagentia beschouwen en dan zou je starten van een reactie
13:39
mengsel waar enkel a2 en b2 aanwezig is en je volgt de opbouw van het reactieproducten a b in functie van de tijd en je ziet dat er AB gevormd wordt.
13:48
Dat betekent dat die reactie doorgaat in voorwaartse richting.
13:59
Maar als je de concentraties van A2 en B2 volgt in functie van de tijd,
14:03
dan zie je dat samen met de opbouw van AB natuurlijk die concentraties omlaag gaan.
14:08
Want je bouwt AB op door A2 en B2 te laten weg reageren,
14:13
maar dat die op het einde van de rit niet op nul terugvallen en dat je dus naar een soort
14:18
reactie mengsel gaat waar zowel die reagentia als die producten nog in aanwezig zijn.
14:22
Opnieuw het feit dat je een soort stationaire evenwicht concentratie bereikt betekent niet dat er in dat reactie mengsel niks meer gebeurt,
14:28
maar dat betekent dat die reactie even vaak per seconde in voorwaartse richting als dat ze in achterwaartse richting doorgaat.
14:37
Je hebt dus opnieuw dat dynamisch evenwicht waarbij je twee flux hebt als het ware de
14:44
reactie in voorwaartse en de reactie in achterwaartse zin die elkaar in balans houden.
14:49
Opnieuw omdat als een minimum van de vrije enthalpie te kunnen zien,
14:55
is het weer essentieel dat dat reactie mengsel dat zijn samenstelling kan aanpassen door de reactie,
14:58
hetzij iets meer in voorwaartse zin, hetzij iets meer in achterwaartse zin te laten verschuiven.
15:05
Als er geen sprake is van zo'n dynamisch evenwicht, is een thermodynamische analyse ook helemaal geen betekenis.
15:11
Opnieuw, die vrije enthalpie wordt geminimaliseerd door het aanpassen van de samenstelling van het reactie mengsel.
15:17
Daar gaat of daar bestaat een wetmatigheid rond die je misschien al wel in het middelbaar gezien hebt en die heet de wet van de massa werking.
15:24
Dat wordt ook wel de Wet van Goldberg en Wagen genoemd en die zegt dat er wat die samenstelling
15:31
betreft dat er altijd een verband bestaat tussen de concentraties van de producten die je vormt.
15:37
En dat is in dit geval A, b en de reagentia die je hebt toegevoegd, dat zijn A2 en B2.
15:43
Je kan natuurlijk die reactie laten starten van verschillende begin concentraties A2 en B2.
15:50
U zou ook kunnen starten met enkel A B aanwezig en kijken naar de omgekeerde reactie.
15:55
Je hebt dus heel veel verschillende mogelijke begin situaties daarvoor.
16:00
Daardoor ook heel veel mogelijke situaties. Maar wat Gilbert en Wagen vaststelden is dat los van die verschillende begin en eind situaties,
16:05
dat die evenwichtstoestand altijd gekenmerkt werd door het feit dat een bepaalde combinatie van concentraties altijd een vast getal heeft.
16:13
En in dit geval, als we de reactie schrijven als a2 plus b2 geeft twee AB,
16:22
dan stelde geen wagen vast dat de concentratie van het product in het kwadraat en het kwadraat komt
16:27
overeen met historische meters gedeeld door de concentratie A2 gedeeld door de concentratie B2,
16:33
dat is altijd een vaste waarde heeft. Dat betekent niet dat die concentraties zelf altijd vast moeten zijn.
16:39
Je hebt speelruimte om die concentraties aan te passen, maar een andere set van concentratie gaat nog altijd aan deze relatie voldoen.
16:46
Dus je kan die drie concentraties niet onafhankelijk van elkaar aanpassen.
16:55
Je kan er misschien twee gaan wijzigen, maar dan gaat het derde zich instellen zodat aan die relatie voldaan is.
16:58
De manier waarop Goldberg en Waag dat hebben uitgelegd, en dat is een wetmatigheid die in het midden van de negentiende eeuw is gepubliceerd,
17:05
is juist vanuit dat idee van een dynamisch evenwicht waar zij argumenteerde van die
17:12
situatie waar die concentraties gecombineerd met het kwadraat en gedeeld door enzovoort,
17:16
een vast getal opleveren, kunnen we eigenlijk begrijpen, zeggen zij uit gelijkheid van die voorwaartse en die achterwaartse reactiesnelheid dat
17:23
er per seconde evenveel omzetting is in voorwaartse zin als in achterwaartse zin.
17:31
Wat ik wil doen is dezelfde wet van de massa werking uitwerken of aantonen,
17:38
maar dan startend vanuit het idee dat dat een thermodynamisch evenwicht is en
17:43
dat het eigenlijk overeenkomt met de samenstelling van laagste vrije enthalpie.
17:47
Los van het idee dat dat een dynamisch evenwicht is of dat je heel specifieke uitdrukkingen.
17:52
Dat is eigenlijk het probleem waar je heel specifieke uitdrukkingen voor dat voorwaarts en achterwaarts tempo moet gaan invoeren.
17:58
Dus chemisch evenwicht is sowieso een dynamisch evenwicht, maar wij kunnen dat gaan analyseren vanuit het idee.
18:06
Dat is de minimalisatie van de vrije enthalpie en niet zozeer die balans tussen een voorwaartse en een achterwaartse reactie.
18:12
Als uitgangspunt nemen om te aan die wetmatigheid van Goldberg en wagen te komen.
18:20
Hoe kunnen we dat doen? Maar als ik het heb over een mengsel van stikstof zuurstof.
18:26
En nu? Waarbij de reactie is stikstof reageert ter vorming van zuurstof.
18:33
En vormt als we de reactie keurig balanceren, tweemaal Anno.
18:41
Ik ga om de ideeën te schetsen starten vanuit een situatie waar je in het begin één en twee hebt en één of twee.
18:50
En ik ga een grootheid invoeren die ik Alpha noem en dat is de omzetting graad.
19:02
Telt eigenlijk hoe vaak? In Mol geteld.
19:14
Die reactie van links naar rechts is doorgegaan. Dus als alfa gelijk is aan nul.
19:19
Betekent dat dat de reactie nog nooit is doorgegaan en dus we hebben enkel.
19:27
Stikstof en zuurstof alfa gelijk aan één betekent.
19:34
De reactie is één mol keer doorgegaan. In dat geval gaan alle stikstof en alle zuurstof verdwenen zijn en heb je twee mol stikstofmonoxide gevormd.
19:39
Één mol die verdwijnt, één mol die verdwijnt geeft je twee mol stikstofmonoxide.
19:50
Je kan een soort figuur maken waar die vordering schraapt.
19:56
Als horizontale as dient. En een voordeel is graad nul betekent enkel stikstof en zuurstof.
20:02
Een vordering tot graad één betekent enkel en nul.
20:10
Dan kan je de vrije enthalpie gaan zien. Als de functie van die vordering geschrapt.
20:15
Herinner je, Vrije enthalpie is een homogene functie, dus we kunnen een vrije enthalpie en ik zal hier een beetje tussenin gaan door krabbelen.
20:23
Altijd schrijven als een som over alle componenten Aantal mol van een component maal de chemische potentiaal van die component.
20:31
Als je niet meer goed weet waar die relatie vandaan komt,
20:40
ga terug naar die slide van twee lessen geleden waarin dat idee van homogene functies invoer.
20:43
Als een functie afhangt van verschillende variabelen, maar er is een subset van variabelen die de eigenschap heeft.
20:48
Als ik elke variabele van die subset maal drie doe, dan gaat de functie waarde maal drie.
20:55
Dan is die functie homogeen in die subset variabelen en kan je ze altijd op die manier gaan schrijven,
21:00
waarbij die chemische potentiaal de partiële afgeleide van g naar een zijn.
21:05
Ik ga dat hier gebruiken, want ik heb. Als de bezettingsgraad alpha is heb ik één min Alpha mol.
21:10
Die stikstof. Ik heb één min al van Mol.
21:18
Die zuurstof en ik heb twee Alfa mol stikstofmonoxide, dus ik kan perfect die vrije enthalpie schrijven in functie van die omzetting graad.
21:23
Dus ik ga hier een verticale as tekenen. Een vrije enthalpie als.
21:36
Stel dat de vrije enthalpie van het stikstof zuurstof mengsel, dan ga je uit.
21:45
Dan ga je uit je thermodynamische tabellen kunnen halen.
21:51
Als ik enkel stikstof en zuurstof heb, dan moet je de vorming vrije enthalpie van stikstof en de vorming vrije enthalpie van zuurstof gebruiken.
21:54
Kan je dus uit thermodynamische tabellen halen. Voor nu kan je dat ook uit die tabellen halen.
22:02
Dan ga je de vergunningsvrij enthalpie van nul gaan gebruiken en heb ik een waarde die bijvoorbeeld daar uitkomt.
22:07
Hoe verloopt die vrije enthalpie tussen die twee uitersten?
22:14
Enkel reagentia alfa nul enkel producten alfa gelijk aan één.
22:19
Als dat de uitdrukking is voor g in functie van alfa, hoe verloopt dan g?
22:24
Als we alfa laten oplopen van nul naar één?
22:31
Die begint hoog, die gaat omlaag en eindigt dan weer hoog. Dus je zegt.
22:45
Ga tekenen. Iets van die jaren.
22:51
Waarom? Maar wacht, jij bent nu de curve aan het construeren.
22:56
Een argument dat je weet waar je moet uitkomen. En dan loop je terug naar waar je vindt dat het eruit moet zien.
23:10
Het juiste. Waarmee ik niet zeg dat het verkeerd is, maar daarmee heb ik nog niet begrepen waarom het er zo uitziet.
23:18
Mijn probleem met jouw uitleg is één en over één en over twee over rechten.
23:24
Dus als ik dat zonder veel nadenken uitzet. Is het zo?
23:35
Waarom ben ik fout en ben jij juist? Zonder.
23:44
Dat is dezelfde argumentatie als ik daarnet hoorde. Als we de vrede geven hem willen is moeten we minimaal we.
24:03
Ja, maar dit kan ook een punt van laagste vrije enthalpie zijn. Waarom is mijn interpretatie van die uitdrukking als een rechte niet juist?
24:08
En dus waarom krijgen we het verloop dat jij zegt? Je ziet.
24:23
Al veel over, over Solidair. Dus geen sprake van.
24:33
Minima. Opnieuw de fameuze slide met de drie vakjes.
24:41
Wat weet je over die chemische potentiaal? Ik wou iets.
24:58
Als dat componenten in de gasfase zijn,
25:09
dan zijn die chemische potentialen beschreven als de standaard waarden die we uit thermodynamische tabellen kunnen plukken,
25:12
plus een bedrag ertoe alleen de partiële druk gedeeld door standaard druk.
25:19
Maar als die reactie door gaat van links naar rechts is er minder en minder stikstof en zuurstof aanwezig en dus
25:25
de partiële druk van die componenten gaat omlaag gaan terwijl er meer en meer stikstofmonoxide aanwezig is.
25:31
De parkeerdruk van die componenten gaat omhoog gaan.
25:38
Dus als we werken met een reactie waarbij die chemische potentialen wijzigen als de samenstelling verandert,
25:42
dan ga je inderdaad in een schema terechtkomen waarbij die vrije enthalpie ergens bij een samenstelling halverwege een minimale waarde gaat bereiken.
25:50
Dat is zeker niet altijd zo. Je kan ook chemische reacties hebben.
26:01
Het is daarom dat het niet zo'n goed idee is om eerst te gaan kijken naar wat moet ik eigenlijk uitkomen
26:06
en dan je redenering aan te passen zodanig dat wat je wenst uit te komen en er effectief ook uitkomt.
26:11
Maar je hebt ook reacties tussen vaste stoffen, bijvoorbeeld barium oxide.
26:16
Kan reageren met titaandioxide ter vorming van barium titaan dat.
26:22
En opnieuw kan je in thermodynamische tabellen naar vormings vrije enthalpie gaan kijken en die gaan bijvoorbeeld daar en daar zitten.
26:31
Maar hier hebben we het over zuivere stoffen.
26:39
Dat betekent dat de activiteit van elk van die stoffen gelijk is aan één en dus in dit geval gaan de chemische
26:42
potentialen helemaal geen functie zijn van de samenstelling en heb je wel degelijk een lineair verloop.
26:49
Tussen beginsituatie en eind situatie en de situatie daar is afhankelijk van de helling die je hier hebt.
26:58
Alles is gewoon reagentia en er wordt niks van producten gevormd of als de helling omgekeerd
27:06
ligt alle reagentia die gaan helemaal verdwijnen en er gaan alleen maar producten overblijven.
27:12
En dus het idee dat we ergens een minimum tussen die twee gaan bereiken is intrinsiek
27:19
gekoppeld aan het feit dat die chemische potentialen wijzigen als de samenstelling verandert.
27:24
En met andere woorden op die manier kan dat systeem de delta geheel gaan aanpassen door zijn samenstelling te veranderen.
27:31
En dan krijg je curves zoals hier getekend. Ik heb ze hier voor stikstofmonoxide, zuurstof en stikstof uitgezet.
27:41
Als je hierop inzoomt op dat stukje vlakbij de Alfa, gelijk aan nul as,
27:48
dan zie je inderdaad een minimum verschijnen en in dit geval gaat dat evenwicht heel sterk aan de enkel,
27:53
stikstof en zuurstof kant liggen en gaat er heel weinig info gevormd worden.
27:58
Maar je ziet wel degelijk een situatie waarbij die alfa verschilt van nul is verschilt van nul op het moment dat de vrije enthalpie het laagst wordt.
28:02
Dus wij gaan op zoek gaan naar die situatie waar die waarde vrije enthalpie niet verandert als we die voor iets graad gaan aanpassen.
28:11
En dat is het punt waar we hier het minimum bereiken en dat is die vordering is graad, die samenstelling waarbij de vrije enthalpie minimaal is.
28:21
Onze aanpak gaat voldoende algemeen zijn, waardoor ook dit soort situaties mee opgenomen gaat worden.
28:32
Je gaat wel zien hoe dat we die gaan kunnen verwerken in de finale analyse.
28:41
Dus waar we naar kijken is verandering van vrije enthalpie door het aanpassen, het wijzigen van de samenstelling.
28:47
Dat betekent, als we het hebben over die DJ term, wat er overblijft.
29:05
De DT term speelt niet mee. De DP term speelt niet mee, want we hebben PT controle.
29:09
De term die toegankelijk is om de vrije enthalpie te verlagen is sommatie over alle componenten mu dni.
29:14
En evenwicht gaat bestaan op het moment dat als wij die samenstelling een klein beetje wijzigen, dat dat geen impact heeft op de vrije enthalpie.
29:24
Met andere woorden de chemische potentialen moeten zodanig zijn dat een kleine wijziging van de samenstelling de vrije enthalpie ongewijzigd laat.
29:33
Ik ga in het algemeen een chemische reactie schrijven als reagentia met een bepaalde stof efficiënt die ik nu hernoem zijn in evenwicht met producten.
29:44
En ik heb een som over verschillende reagentia en ik heb ook een som over verschillende producten.
29:58
Dus grote en grote P duiden verschillende reagentia en producten aan die Griekse letter.
30:05
Nu met de index r. Dat is de sleutel om efficiënt van een bepaald reagens.
30:12
Het is een zeer algemene manier om chemische reacties te schrijven.
30:19
Bijvoorbeeld het verdamping licht waar we het daarnet over hadden kan je ook op die
30:23
manier schrijven als water in de vloeistof fase is in evenwicht met water in de gasfase.
30:27
Één reagens, één product en de stokjes zijn allebei één.
30:36
Dus wat we vorige les over fysisch evenwicht gezegd hebben, gaan we eigenlijk nu veralgemenen.
30:41
De relaties die we gaan opstellen gaan ook van toepassing zijn op dat fysisch evenwicht.
30:45
Als ik een onderscheid maakt tussen reagentia en producten, ga ik die dg gaan schrijven als sommatie over alle reagentia.
30:54
Chemische potentiaal van een reagens, verandering van het aantal mol van dat reagens plus sommatie over alle producten
31:03
Chemische potentiaal van een product Mol Verandering van het aantal mol van dat product.
31:12
Waarom maak je het onderscheid tussen reagentia en producten? Omdat ik.
31:19
Ik ga al die definities van producten en reagentia proberen schrijven in termen van één enkele variabelen.
31:24
En ik ga die variabelen. Mijn lievelings letter.
31:32
De Griekse letter C moet je oefenen om die te leren schrijven.
31:36
De letter is de. En dat noemen we de vordering graat van de redactie.
31:42
Is een beetje analoog aan die omzetting gratis die ik daarnet heb ingevoerd.
31:54
Maar telt in Mol het aantal keer dat de reactie van links naar rechts is doorgegaan.
31:59
Dat betekent dat ik een verandering van het aantal mol van een bepaald product kan
32:07
schrijven als verandering van de vordering graad maal die strook van meters efficiënt.
32:13
Als de reactie een aantal keer drie x extra doorgaat van links naar rechts, dan maak ik een aantal mol van een product bij.
32:22
Gelijk aan dus toch efficiënt van dat product.
32:33
Maal de verandering van vorderingen graad, want per keer dat ze doorgaat vorm ik nu één mol van dat product.
32:36
Analoog kan je een verandering van het aantal mol van een reagens dat gaat verdwijnen als de reactie extra doorgaat.
32:44
Van links naar rechts kan ik schrijven als min strookje metrische coëfficiënt van dat reagens maal de verandering van die vordering graad.
32:52
En dus die DJ en rietjes die daar staan en die ideetjes die daar staan kan je eigenlijk
33:02
allemaal samenvatten in termen van één enkele dictie een verandering van vordering schraapt.
33:08
En de cd. Laat zich op die manier schrijven als ik ga producten links gaan zetten, als ik de DNP substitueren, dan krijg je de Griekse letter.
33:14
Nu, dat is de stoter metrische coëfficiënt Griekse letter nu van dat product.
33:25
Dat is de chemische potentiaal vaardigheid.
33:30
In het Griekse alfabet is handig. Min sommatie over alle reagentia strooi met deze chemische potentiaal.
33:33
En dan het summum van het Griekse alfabet.
33:44
Verandering van verordening schraapt die dictie en bij evenwicht moeten die chemische potentialen zodanig zijn dat dat zaakje gelijk is aan nul,
33:46
dat een verandering van vordering schraapt.
33:58
Het klein beetje extra laten doorgaan van de reactie van links naar rechts dat dat geen impact heeft op de vrije enthalpie,
34:00
dat is dat minimum van die curve waar we naar op zoek zijn.
34:07
Je ziet dit gaat nul zijn voor eender welke variatie van actie, op voorwaarde dat of wanneer die voor factor gelijk is aan nul.
34:13
Dus we kunnen die advies herschrijven als sommatie over alle producten.
34:24
Meters efficiënt chemische potentiaal min sommatie over alle reagentia cytochroom
34:33
secco efficiënt chemische potentiaal Dat moet bij evenwicht gelijk zijn aan nul.
34:41
Wat gaat er gebeuren als die combinatie negatief is?
34:56
Als die sommatie over alle producten nu PMP informatie over alle reagentia stoten minder efficiënte chemische potentiaal.
35:02
Wat gaat er gebeuren als die combinatie negatief is?
35:11
Of dat negatief is, dan zie je dat de vrije enthalpie omlaag kan gaan voor een toename van vordering gaat.
35:16
Zolang dat die combinatie negatief is, gaat die reactie doorgaan in voorwaartse zin.
35:26
Dus je kan hier ook een kleiner dan teken bijzetten gelijk aan nul betekent.
35:33
De chemische potentialen zijn zodanig dat betekent de samenstelling is zodanig dat er evenwicht is.
35:39
Als dat linkerlid kleiner is dan nul, dan is die samenstelling zodanig dat die reactie,
35:46
als ze doorgaat in voorwaartse zin, de vrije enthalpie nog verder kan verlagen.
35:52
Als die combinatie positief is, dan kan de DG omlaag gaan voor een negatieve waarde van actie.
35:59
Dat betekent dat dat evenwicht of die reactie in achterwaartse zin gaat doorgaan.
36:06
Dus het vergelijken van die combinatie van stoot en chemische potentialen met nul leert ons of een reactie in evenwicht is,
36:11
dan wel of ze in voorwaartse of in achterwaartse zin gaat doorgaan.
36:22
Kleiner dan nul voorwaarts groter dan nul achterwaarts gelijk aan nul evenwicht.
36:27
Heeft te maken met als we die curve hier maar tekende. Heeft te maken met die hellingen.
36:33
Nul betekent DG op de IS nul.
36:39
We zitten in een minimum positief. Dan gaan we in achterwaartse zin negatief.
36:43
Dan gaan we in voorwaartse zin naar dat minimum toegaan bewegen.
36:49
Wie kan er betekenis geven aan die combinatie van Griekse letters hier?
37:03
Stel dat bij een gegeven samenstelling. Die reactie.
37:17
Één keer doorgaat van links naar rechts.
37:22
Als ze eenmaal door gaat van links naar rechts,
37:27
dan vorm je nu p mol van elk product bij en de verandering van vrije enthalpie is dan per product dat er bijkomt.
37:30
De combinatie stootte met centraal chemische potentiaal.
37:39
Dat is de toename van vrije enthalpie als je extra product creëert.
37:44
Minder afname van vrije enthalpie als reagens verdwijnt wanneer die reactie één keer doorgaat van links naar rechts.
37:49
Dat is dus de reactie vrije enthalpie. De vrije enthalpie die gepaard gaat met de producten.
37:59
Minder vrijheid die gepaard gaat met de reagentia.
38:05
Als de reactie eenmaal een keer doorgaat niet bij standaard omstandigheden, maar bij de feitelijke samenstelling van het reactie mengsel.
38:09
Om die mu P's vorm te geven moet je daar bijvoorbeeld de feitelijke partiële drukken van al die gas componenten gaan invullen.
38:17
Nu Wat hebben we gezegd over chemische potentialen, zelfs als we ze niet kennen?
38:25
We gaan die altijd in twee stukken trekken en we hebben aan de ene kant.
38:29
De standaard chemische potentiaal die we uit tabellen kunnen aflezen.
38:36
En als ik die chemische volgens jou van het product splits in twee stukken standaard waarde plus RTL en
38:41
activiteit kan ik de standaard waarden in het linkerlid laten staan en ik ga dat ook doen voor de reagentia.
38:47
Ik laat daar ook de standaard waarden in het linkerlid staan.
38:55
En wat ga ik naar het rechterlid brengen? Dat zijn alle combinaties die met activiteiten te maken hebben.
38:59
Daar krijg je dus. Min rte sommatie over alle producten logaritme van een bepaald product tot een exponent gelijk aan de stock sector efficiënt.
39:06
In principe staat er nu p maal de logaritme, dat is de logaritme van het argument tot die exponent, tot een factor die er voor staat.
39:21
En ik heb een bedrag min rte analoog voor de reagentia activiteit van die reagentia tot hun stock efficiënt.
39:30
En wat gebruik ik? Dat is de algemene opsplitsing van een chemische potentiaal in een standaard waarde plus een bedrag artikelen activiteit.
39:41
Dat was het derde vakje in die slide die je uit je hoofd moet kennen.
39:50
Op dit moment kennen we de activiteiten in twee situaties zuivere stoffen en componenten in een ideaal gasmengsel.
39:57
We gaan later vandaag activiteiten leren kennen.
40:04
Ook voor componenten in oplossing waardoor we onze chemische evenwichts relatie op meer situaties gaan kunnen toepassen.
40:07
Het kan zijn dat ik er pas volgende week aan begin, maar dat gaat een volgende zijn die we leren kennen.
40:15
Wat ik hier doe is die relatie in twee stukken halen.
40:21
Dit is iets dat enkel afhangt. Van de standaardwaarden.
40:27
Dus daarvoor ga je naar een tabel en haal je getallen uit.
40:33
Een tabel is dus onafhankelijk van de feitelijke concentraties van reagentia of producten waarmee je gestart bent.
40:36
Dat is puur een eigenschap van de reactie waar het over gaat en de omstandigheden verstaan de temperatuur waarbij die reactie gebeurt.
40:44
Alle afhankelijkheid van de samenstelling zit daar.
40:53
En nog altijd kan ik met een kleiner dan en grotere dan teken meepakken.
40:59
Kleiner dan sorry gelijk betekent evenwicht kleiner, dan betekent reactie in voorwaartse zin groter, dan betekent de reactie in achterwaartse zin.
41:04
Wie heeft er een naam voor wat er hier staat? Dat is de vormings vrije enthalpie van de producten.
41:16
Eenmaal de stof gemeten c.q. efficiënt minder vorming vrije enthalpie van de reagentia.
41:28
Maal die meters. Wat is dat? Wie kan daar een naam aan geven?
41:35
We hebben daar een beetje tijd in gestoken om dat te benoemen en dat uit te leggen.
41:50
Ik dacht dat je een beweging maakte om. Ik dacht dat je een beweging maakt door te zeggen.
42:03
Volgens velen een kopie van het product basisdocument als een overwegend vergunningsvrij interactief reagens basisdocument.
42:11
Als ik vind, wat bereken je dan? Zal ik voor een bepaalde reactie.
42:17
Ik zal een gas reactie gebruiken. Stikstof en waterstof reageren ter vorming van ammoniak.
42:33
Stel dat ik je vraag om voor die reactie de reactie vrije enthalpie te berekenen.
42:43
Hoe doe je dat? Naar welke getallen ga je op zoek?
42:48
Dat is de Delta G. Ja, en hoe doe je? Hoe bereken je dat? Het staat niet meer open.
43:03
Misschien. Ik ga even kijken. Nee, het staat niet meer open.
43:36
Even geduld. Kijk eens hier.
43:43
Wacht, wacht, wacht, wacht, wacht. De slides van les dertien zijn eigenlijk wel belangrijke slides, dus ik heb het al.
44:09
Ik heb al benadrukt die slides, die drie vakjes, dat je die uit je hoofd moet kennen.
44:25
Maar er zijn meerdere dingen die je daarvoor moet begrijpen. Standaard vorming vrij standaard reactie Vrije enthalpie sommatie over alle producten.
44:28
In eerste instantie hadden we de absolute vrije enthalpie die niet bestaat,
44:40
maar je weet dat je dat door de vorming vrije enthalpie mag veranderen die je uit thermodynamische tabellen haalt.
44:44
Maal gesommeerd om voor alle producten min hetzelfde voor de reagentia maal de gemiddelde coëfficiënten gesommeerd over alle reagentia.
44:49
Dat is exact wat er hier staat. Stel je weet dat ik ooit vier maal vergunningsvrij enthalpie van een product stook.
44:58
Efficiënt, maar volledig vrij van elke reagens. Dus hier hebben we de standaard reactie vrije enthalpie.
45:07
Die we zien verschijnen uit de combinatie van die standaard waarden.
45:18
Dit is wat je vindt voor elk product en voor elk reagens.
45:22
Intern dynamische tabellen onder het kolommetje Delta G.
45:26
Vorming. Ik heb kleiner dan gelijk aan of groter dan.
45:30
We gaan de Minnertsga. Naar het andere lid brengen.
45:41
Dat is de RTE die daar staat. Wat krijgen is een sommatie van logaritme.
45:47
Dat is ook het logaritme van het product van die activiteiten tot hun stock gemeten,
45:53
efficiënt gedeeld door het product van die activiteiten over alle reagentia tot hun stokje efficiënt.
46:00
Typisch gaat men links en rechts de exponentiële nemen en krijg je de exponentiële van de delta G nul gedeeld door R2.
46:21
Die bij evenwicht gelijk is, maar die ook kleiner dan of groter dan kan zijn.
46:34
Omdat ik links en rechts maal -1 gedaan heb,
46:40
gaat de betekenis van die kleiner dan één groter dan tekens wisselen en rechts krijg je het product
46:44
over alle producten activiteit tot de socialistische coëfficiënt en hetzelfde voor de reagentia.
46:52
Opnieuw. Wat er links staat is een eigenschap van de redactie waar we het over hebben bij een bepaalde temperatuur.
47:02
Het enige wat je moet weten is wat zijn de reagentia, Wat zijn de producten?
47:10
Je haalt Delta C0 uit een thermodynamische tabel en natuurlijk moet je werken bij de temperaturen waarin je geïnteresseerd bent.
47:14
De temperatuur waarover het gaat. Dit heet daarom de evenwichts constante.
47:21
Omdat het voor een gegeven reactie bij een gegeven temperatuur is dat een vast getal.
47:29
Dit is anders. Dit hangt af van de activiteiten.
47:36
Denk als het over een GAS reactie gaat aan de partiële drukken.
47:40
Dat gaat dus afhangen van de samenstelling. Als de reactie in voorwaartse of in achterwaartse zin evolueert.
47:45
Die is constant en die gaat niet aangepast worden. Wat er wel gaat veranderen is het rechterlid.
47:53
Een reactie gaat altijd die samenstelling aannemen waarbij het rechterlid gelijk is aan het linkerlid.
48:00
Dit heet het reactie quotiënt. Coachend verwijst naar die verhouding.
48:07
Als het reactie quotiënt kleiner is dan de evenwichts constanten, dan gaat die reactie verder gaan.
48:23
In voorwaartse zin is het reactie quotiënt groter ga je in achterwaartse zin.
48:28
Evenwicht bereik je als de twee aan elkaar gelijk zijn. Opnieuw, je moet die vergelijking, die evenwicht vergelijking, goed leren interpreteren.
48:34
Die evenwicht constante is een vast getal gegeven de reactie gegeven.
48:43
De temperatuur is een vast getal. Voor elke samenstelling van het reactie mengsel kan je haar reactie quotiënt berekenen.
48:48
Denk aan een gas reactie. Van zodra je de partiële drukken hebt van alle reagentia en alle producten kan je dan reactie quotiënt berekenen.
48:56
Uit de vergelijking tussen reactie quotiënt en evenwicht constante kan je begrijpen kan je inzien wat er gaat gebeuren.
49:04
Zolang het reactie quotiënt kleiner is dan de evenwicht constante, gaat de reactie doorgaan in voorwaartse zin.
49:12
Als de reactie quotiënt groter is dan de evenwicht constante, ga je in achterwaartse zin doorgaan.
49:18
De reactie bereikt evenwicht op het moment dat de reactie quotiënt en de evenwicht
49:24
constanten aan elkaar gelijk zijn en dat is de evenwichts voorwaarde die we gaan gebruiken.
49:28
Laat ons eens kijken wat het wordt als we terugkeren naar dat verdamping en condensatie evenwicht.
49:36
Ik heb daarnet proberen zeggen de aanpak is algemener en bevat ook wat we vorige week gezien hebben.
49:43
Althans als het reagens bijvoorbeeld vloeibaar water is, gaat dit nog wat omhoog zitten.
49:52
Als het reagens vloeibaar water is en het product is gasvormig water,
50:02
dan hebben we nu geleerd evenwicht gaat bestaan op het moment dat de exponentiële van min
50:09
delta geen nul gedeeld door R2 gelijk is aan het reactie quotiënt het reactie quotiënt.
50:14
De naam zegt het zelf is een breuk. Wat krijg je in de teller?
50:21
Dan moet je gaan kijken naar je producten. En de activiteit van elk van die producten moeten we nemen tot zijn stokje efficiënt.
50:26
Wat komt er in de noemer? Dan moeten we kijken naar de reagentia Activiteit wordt genomen tot zijn structuur efficiënt.
50:34
Producten. Hier is water in de gasfase.
50:42
Wat is de activiteit van zo'n component in de gasfase? Hier is de fameuze slide.
50:46
Wat is de activiteit van een competente gastvrouw? Dat is de partiële druk van water in de gasfase te delen door standaard druk.
51:12
Dat is de teller die we hebben uitgewerkt. Wat is de noemer activiteit van in dit geval het reagens zuiver vloeibaar water?
51:22
Is één. En je ziet dat je meteen komt in het schema dat we daarnet hadden.
51:30
De dampdruk is te berekenen als standaard druk maal exponentieel van min delta,
51:34
geen nul gedeeld door de delta, geen nul is hier opnieuw die verdamping.
51:39
Vrije enthalpie standaard waarde voor de gasfase min standaard voor de vloeistof fase.
51:44
Dus wat we vorige les gedaan hebben zit mee ingepakt, zit mee verpakt in die chemische evenwichten relatie die we nu hebben uitgewerkt.
51:50
En van zodra je een omzetting kan schrijven in die chemische reactie terminologie
51:58
kan je die evenwichts relatie zoals we daar hebben uitgewerkt gebruiken.
52:05
Evenwicht constante moet gelijk zijn aan het reactie quotiënt bij evenwicht.
52:10
We gaan dat vandaag minstens. We gaan zien hoe ver we geraken.
52:19
Minstens proberen toepassen op gast reacties.
52:23
Dat is leuk om oefeningen over te geven, maar het helpt je ook een beetje hoe je met die evenwichts relaties aan de slag moet gaan.
52:27
Dus wat hebben we hier allemaal? Even kijken. Dus in de sla op de slides.
52:44
Dat is de theorie die ik ook op bord heb gebracht.
52:52
Invoeren van de vorderingen graad De deelnemers schrijven in termen van een discussie de vrije enthalpie verandering schrijven voor factor Dixie.
52:54
Verschil tussen voorwaarts, achterwaarts en evenwicht. Die vormfactor moet gelijk zijn aan hulp bij evenwicht.
53:03
Het invoeren van de opsplitsing standaardwaarden die we in tabellen vinden en al de rest die we RTL activiteit noemen.
53:09
Alle samenstellingen afhankelijkheid zit in die term.
53:18
RTL inactiviteit laat ons toe de event voorwaarden te splitsen in twee delen een deel dat de evenwicht constanten
53:21
gaat opleveren en afhangt van die delta geen nul en het ander deel waar al die activiteiten samenkomen.
53:28
En dat gaat afhangen. Dat is het reactie quotiënt.
53:34
Dat is het product van al die activiteiten dat is gemeten, zowel voor de producten als voor de reagentia.
53:36
En dan heb je. De interpretatie daarvan en waarbij we gaan kijken naar het verband tussen het reactie quotiënt en de evenwicht.
53:43
Constanten kleiner dan groter dan of gelijk aan waar het gelijk aan teken betekent.
53:53
Dat is die samenstelling waar er evenwicht bestaat. Groter dan of kleiner dan is.
53:58
Voorwaartse zin of achterwaartse zin. Wat krijgen we als we.
54:02
Die algemene relatie toepassen. Op een gas reactie.
54:23
En een voorbeeld van een gas reactie is wat we daarnet hadden. Het stikstof reageert met zuurstof.
54:31
De twee belangrijkste samenstellingen in de lucht ter vorming van stikstofmonoxide Een
54:36
vervelende reactie omdat die kan optreden in verbrandingsmotoren bij hoge temperatuur.
54:42
En die stikstofoxiden geven aanleiding tot wat heet.
54:46
Fotografische smog is dus een vervelende reactie, maar ze kan wel optreden als we daar die relaties op toepassen.
54:49
Wat krijgen we? Wat worden de activiteiten van elk van die producten en reagentia?
54:55
Als dat componenten in een ideale gasfase zijn, dan gaan we als reactie quotiënt hebben.
55:01
Partiële druk van Enno in de gasfase gedeeld door standaard druk tot zijn efficiënt betekent
55:07
in het kwadraat gedeeld door partiële druk van zuurstof in één van standaard druk.
55:16
Partiële druk van stikstof in eenheden van standaard druk en dat moet bij evenwicht gelijk zijn aan de evenwicht constante.
55:23
Ik ga reviews constanten hier altijd noteren als k met twee.
55:36
Toevoegingen een subscript p. En dan in dit geval relevant dat bolletje dat erboven staat.
55:43
Het subscript p dat ik hier gebruik verwijst naar het activiteit model dat er bijhoort.
55:50
Dat betekent dat activiteiten in partieel geschreven worden een nulletje betekent.
55:55
Die evenwicht constanten gaan we rechtstreeks uit een delta geen nul het gestabiliseerde waarden gaan afleiden.
56:01
Waarom is die notatie belangrijk? Omdat partiële drukken niet altijd de meest handige grootheid zijn om evenwichten te gaan beschrijven.
56:11
Vaak is het interessanter om samenstellingen te beschrijven in termen van mol fracties.
56:24
Waarbij de Mol fractie van een component A. Geschreven als x a.
56:33
Het aantal mol van die component is gedeeld door het totaal aantal mol in de gasfase.
56:39
Fracties geven een beter beeld over samenstelling mol fracties.
56:47
Nulpunt één betekent 10% van de moleculen is a mol fractie nulpunt vijf betekent 50% van de moleculen is A.
56:50
Hoe kan je van een uitdrukking in partiële drukken naar een uitdrukking een mol fracties overgaan?
57:03
Partiële druk van A. Kan je altijd schrijven als de totale druk die er is maal de mol fractie van die
57:11
component is de bijdrage van die component tot de totale druk van dat ideaal gas.
57:21
Dus als je in totaal een evenwicht hebt zoals hier.
57:30
Producten over product over alle producten Partiële druk tot diastolische Metrische.
57:34
Coachend. Alle reagentia partiële druk tot histogram risico efficiënt is de evenwicht constante KP nul.
57:39
Dan kan je overgaan naar een reactie coachend.
57:50
Niet in partiële drukken, maar in fracties.
57:54
Door deze substitutie. Je ziet dat je dat ook kan schrijven als.
57:57
Product. Partiële druk die ik substitueren reken Mol fractie.
58:03
Tot die social media zijn efficiënt. Je krijgt ook een totaal druk.
58:11
Die ga ik seffens meepakken. Per mol fractie die gesubstitueerd krijg je een mol fractie niet A maar B van een reagens tot zijn stokje tussen.
58:16
Proficiat. Prostitutie die ik doorvoer.
58:26
Ik heb totaal tot nu pijn in de teller en ik kreeg heb totaal tot nu er in de noemer.
58:31
Netto effect gaat een pet totaal zijn.
58:38
Gedeeld door P0 ontstaat eigenlijk partiële druk gedeeld door standaard druk partiële druk gedeeld door standaard druk.
58:42
Ik kan die gedeeld of standaard druk mee pakken in die p totaal en dat gaat een wat we noemen delta nu zijn.
58:50
Dat is het verschil histogram. Als ik efficiënt aan de product kant min de som van de stock koffietenten aan de reagens kant en je ziet
58:58
dat dat voortkomt uit die substituties en dat gaat bij evenwicht opnieuw gelijk zijn aan die KP nul.
59:06
Heb ik nog ergens een plekje om verder te werken? Jawel, ik ga hier verder werken.
59:20
Ik zal mijn relatie opnieuw opschrijven. We hebben product over alle producten XP nu p x er nu er.
59:25
En dan moet gelijk zijn aan de KPN nul. Maal die combinatie.
59:35
P Totaal. In eenheden van standaard druk en als ik hem naar het andere lid breng dan krijg ik daar een min dat dan nu in plaats van een plus delta nu.
59:42
Dit kunnen we aanduiden als een reactie coachend. Maar geschreven in Mol fracties.
59:57
Het is eigenlijk wat we hier hebben in termen van activiteiten, maar elke activiteit vullen we nu in als een mol fractie.
1:00:08
Is een reactie quotiënt. De index x wijst naar het feit dat we als activiteit de mol fracties gebruiken en
1:00:17
dan zie dat dat bij evenwicht gelijk moet zijn aan de oorspronkelijke KP nul.
1:00:25
En de KPI nul. Dat is die exponentiële min.
1:00:31
Delta G nul gedeeld door RTE. Die haal je dus uit de tabel hier de delta G waarden.
1:00:34
Vandaar dat ik er dat nulletje bij zet, maar is niet gewoon gelijk aan de KPI.
1:00:40
Nul is gelijk aan de KPI nul maal pi totaal gedeeld door P nul tot de min delta nu.
1:00:45
Die combinatie ga je vaak aangeduid zien worden als opnieuw een evenwicht constanten, maar dan in een mol fractie model.
1:00:54
Ik zet hier het nulletje er niet bij omdat die enige constante die ga je niet rechtstreeks uit een delta geen nul halen.
1:01:04
Om die constante te bepalen moet je de KPI nul bepalen.
1:01:11
Die komt rechtstreeks uit de delta. Geen nul één combineren met TP totaal gedeeld door P nul tot de min delta nu.
1:01:15
Dus dat is een erg constant die wel kan verschillen van die KPN nul afhankelijk van de waarde van delta nu.
1:01:22
En je kunt je. Je kan je evenwicht dus behandelen, hetzij door een QP een reactie quotiënt in partiële drukken gelijk te stellen aan een KPI nul.
1:01:30
Dat is een manier om het te doen. Exact hetzelfde evenwicht ga je berekenen door een x gelijk te stellen aan een X.
1:01:42
Afhankelijk van de vraagstelling is de ene route of de andere route het eenvoudigst
1:01:51
en doorgaans zijn de vragen zo geformuleerd dat die route de eenvoudigste is.
1:01:56
Waar je bij moet opletten is dat je natuurlijk altijd de manier van schrijven van van je
1:02:02
reactie coachend combineert met de juiste manier van berekenen van de evenwicht constanten.
1:02:07
Je kan kiezen voor een mol fractie model. Je kan kiezen voor een partieel model, maar je moet natuurlijk consistent zijn.
1:02:13
Als je reactie coachend in mol fractie schrijft, moet je de daarbijhorende evenwicht constanten gebruiken.
1:02:19
Als je je reactie coachend in partieel gebruikt, moet je daarbijhorende evenwicht constanten gebruiken.
1:02:25
Laat ons eens proberen bekijken hoe we dat in de praktijk kunnen aanpakken dan.
1:02:33
Ik zal in eerste instantie verder gaan met die reactie tussen stikstofmonoxide, tussen stikstof en zuurstof ter vorming van stikstofmonoxide.
1:02:41
Dus we hebben een twee. En O twee geven ons twee maal nul.
1:03:03
Als ik in een fractie model werk. Heb ik een evenwicht constante kicks.
1:03:20
Die gelijk moet zijn aan de reactie quotiënt, ook in zo'n mol fractie model.
1:03:27
Wat gaat dat worden, dat reactie quotiënt? Wie kan de teller en de noemer invullen?
1:03:32
De algemene relatie voor een reactie Coachend hier staat en we gebruiken als mol fracties, wat krijgen we dan?
1:03:49
In de teller moeten we gaan kijken naar wat zijn de producten?
1:04:04
De activiteit van elk product, in dit geval dus de Mol fractie, moet daar komen, telkens tot een macht gelijk aan zijn stokje.
1:04:07
Efficiënt dus. Welke activiteiten gaan daar verschijnen en tot welke macht?
1:04:14
Waarom? Wat bedoel je met één vierde en één tweede? De vraag is zijn op dit moment onbekenden die zouden moeten bepalen uit de evenwichts relatie?
1:04:36
Het enige dat we kunnen zeggen is de fractie van. En hoe.
1:04:48
Al twee. En twee moeten zodanig zijn.
1:04:55
Dat de Mol fractie van één kwadraat gedeeld door de Mol fractie van CO2 gedeeld door
1:05:02
de Mol fractie van N2 gelijk is aan die even constante als deze is voorwaarde.
1:05:08
We hebben een vergelijking. Voorlopig in drie onbekende.
1:05:19
Dan gaan we niet oplossen. Het tweede dat we kunnen invoeren is.
1:05:25
Als ik start met. Voor O twee en twee en O als ik één mol heb één mol en niets bij het begin en ik
1:05:34
ga hier de omzetting schraap schrijven of de dekkingsgraad schrijven als nul.
1:05:50
Wat ik aan het doen ben is een verband zoeken tussen die verschillende mol fracties, omdat die reactie natuurlijk een vaste structuur met drie heeft.
1:05:56
Als die gaan starten met nul mol en nul mol.
1:06:06
Om het algemener te maken als de reactie alfa en nul keer is doorgegaan.
1:06:09
Dus ik schrijf mijn vorderingen graad als evenredig met het aantal mol waarmee ik start.
1:06:17
Dan is er een aantal mol zuurstof verdwenen per mol.
1:06:25
Keer dat de reactie doorgaat verlies ik ook één mol zuurstof.
1:06:28
Dus als die reactie alfa en nul keer doorgaat dan ben ik min alfa en nul mol zuurstof kwijt.
1:06:32
Ik ga ook min en nul mol. Stikstof verliezen en ik kreeg twee alfa en nul mol.
1:06:39
In de plaats. Wat heb ik als totaal aantal.
1:06:49
Ik heb twee en nul. Wat komt er bij?
1:06:56
Alfa? Alfa plus twee Alfa. Er verandert niets.
1:07:00
Totaal aantal Mol blijft gelijk. Je verliest twee mol en je krijgt twee mol in de plaats.
1:07:05
Met andere woorden de delta nu van deze reactie is nul.
1:07:11
Som van de stock gemiddelde koffie centen rechts en links is in beide gevallen is twee.
1:07:16
Hoeveel geeft dat in totaal? Dat is één mijn Alfa.
1:07:23
En nu? Je hebt één alfa en nul en ik heb twee alfa en nul en nul en mijn totaal aantal mol is twee en nul.
1:07:29
Wat krijg je dus als Mol fracties? Dat is het aantal mol van een component gedeeld door het totale aantal mol.
1:07:41
Geef mij één mijn alfa. Één nul gedeeld door twee en nul is één min van gedeeld door twee.
1:07:50
Dat is de Mol fractie voor de stikstof. Één en al van gedeeld door twee.
1:07:58
Dat is ook de Mol fractie voor Sorry Diva die zuurstof en stikstof en dan twee
1:08:03
alfa en nul gedeeld door twee en nul één keer alfa dat is de Mol fractie en.
1:08:10
Dus ik heb een vergelijking inzicht vergelijking in drie onbekende mol fracties via een massa balans heet dat.
1:08:16
Kan ik al die mol fracties in één enkele variabele schrijven.
1:08:25
Dat is die voor de ruggengraat waar we naar op zoek gaan, waarbij er evenwicht gaat bestaan.
1:08:31
Als ik die drie mol fracties invult invul krijg ik X is gelijk aan mol fractie en dat is alfa kwadraat mol fractie CO2 dat is één
1:08:36
min Alfa mol fractie N2 is één min alfa en ik had een gedeeld door twee en een gedeeld door twee geeft mij vier alfa kwadraat.
1:08:49
En zo bouw je die evenwichts vergelijking die in drie onbekenden geschreven wordt.
1:08:59
Om naar een vergelijking in één enkele onbekende.
1:09:05
Hieruit los je. Alfa op.
1:09:11
Dat geeft u de omzetting,
1:09:16
schraapt bij evenwicht en dan gebruik je opnieuw je formules die alfa koppelen aan een mol fractie om de samenstelling te bepalen.
1:09:17
En dat is de manier van werken om zo'n chemie stevig in de gasfase op te lossen.
1:09:26
We starten van een uitdrukking voor de evenwichts vergelijking, voor de chemie, voor de reactievergelijking.
1:09:32
Sorry, dat vertaal je in de vergelijking reactie quotiënt.
1:09:36
Aan de ene kant evenwicht constant aan de andere kant.
1:09:42
Je kan kiezen voor het activiteit model.
1:09:47
Doorgaans is het een goede keuze om te werken met een mol fractie model, maar je gaat dezelfde relatie natuurlijk uitkomen.
1:09:50
Of je gaat zelf de evenwichtssituatie berekenen door met een partieel druk model te werken.
1:09:57
Keuze van het model kan nooit de ligging zelf veranderen.
1:10:02
Als we kiezen voor een model fractie model. Het voordeel daarvan is dat die massa balans zich heel gemakkelijk laat uitwerken.
1:10:07
Je kijkt wat heb ik in het begin bijvoorbeeld enkel zuurstof en stikstof, geen enzo.
1:10:15
Als de vordering is graad alfa en nul bedraagt. Hoeveel stikstof en zuurstof ben ik kwijt?
1:10:20
Hoeveel en nul heb ik gevormd? Je kan de aantal mol schrijven in termen van alpha en dus schrijf je de mol fracties in termen van alpha.
1:10:26
Wacht nu ook niet mee. Ah nee, nee, nee, nee, die moet daar weg hè?
1:10:39
Het is een fractie Stikstof is één en nul gedeeld door twee nul.
1:10:50
Dank je. En goed.
1:10:54
In dit geval krijg je een kwadratische vergelijking waar je alpha het oplost.
1:10:58
Om de X te berekenen. Wat heb je daarvoor nodig?
1:11:18
Een aantal zaken. Je moet het totaal druk kennen.
1:11:35
Dat moet gegeven zijn. Je zoekt een evenwicht bij een bepaalde totale druk.
1:11:39
Dan moet gegeven zijn de delta Nu die haal je uit de evenwichts vergelijking en dat is de som
1:11:43
van de sterke efficiënte aan de product kant en de som van de centen aan de reagens kant.
1:11:49
In dit geval is dat nul en valt die term dus weg en is de X gelijk aan de KP nul, de KP nul zelf.
1:11:55
Daarvoor ga je naar je thermodynamische tabel en heb je de delta geen nul die je moet bepalen gedeeld door R2.
1:12:02
Opnieuw. Voor die delta geen nul kan je het principe gebruiken.
1:12:10
Hij is zelf altijd afhankelijk van temperatuur, maar als je bij een andere temperatuur werkt als waarbij de data van je tabel gegeven zijn,
1:12:13
kan je typisch delta en Delta als temperatuurs onafhankelijk beschouwen en bereken je op die manier een benaderende delta geen nul.
1:12:22
Dus als je tabellen hebt bij twee 98 Kelvin en je moet een evenwicht berekenen bij 500 Kelvin, gebruik je delta geen nullen uit de tabel,
1:12:30
maar gebruiken Delta H0 en je delta tabel en combineer ze om een Delta GL bij 5000 Kelvin te bepalen.
1:12:38
Ik heb die zaken ook in de slides gezet, dus de stappen die je nodig hebt om zo'n evenwicht uit te werken.
1:12:47
Je start van je reactievergelijking, je schrijft je evenwichts vergelijking uit.
1:12:54
Bestaat het? Een reactie uit een evenwicht constante en een reactie quotiënt uit reactie quotiënt.
1:12:59
Daar zitten al die verschillende activiteiten in en die breng je samen.
1:13:05
Al die reacties, al die mooie reacties breng je samen aan de hand van die massa balans.
1:13:14
Alles wordt geschreven in één enkele variabele die je de omzetting schrijft kan noemen.
1:13:19
Die haal je uit de evenwichts relatie.
1:13:24
Daarmee ga je terug naar de massa balans om de verschillende mol fracties te bepalen om de evenwicht constanten eruit te halen.
1:13:26
Pas op, je kan niet enkel en alleen met je delta geen nul werken als je per se als je
1:13:33
totale druk verschilt van standaard druk en je delta nu is niet gelijk aan nul,
1:13:38
is je X verschillend van je KP nul. En ga je dus voor andere, drukke, andere levens constant te gaan bepalen.
1:13:43
Dus de dat geen nul. Dat staat hier nog eventjes uitgelegd. Daar hebben we vorige les aldaar al al aangehaald dat je naar die delta moet gaan kijken,
1:13:55
dat deze typisch als temperatuurs onafhankelijk moet beschouwen. Laat ons er eens mee aan de slag gaan.
1:14:03
Ik zie dat bijna tijd is om pauze te nemen, maar ik wil toch eventjes.
1:14:09
Een voorbeeld uitwerken waarbij we gebruik maken van de tabellen boekje en ik ga.
1:14:17
Kijk hier naar een reactie waarbij we O twee in evenwicht bekijken met O drie.
1:14:37
Waarbij we toch een koffietentje kunnen kiezen. En ik ga hier werken met drie en halve zuurstof geeft één op drie.
1:14:49
Misschien toch beter om eerst 5 minuten pauze te nemen en dan met deze strijd verder te gaan.
1:14:57
Ook ik stel voor dat we verder doorgaan. Oké.
1:30:33
Dus we gaan die. Aanpak van gas evenwichten is toepassen op.
1:31:06
Die vorming van ozon uit die zuurstof.
1:31:15
We gaan in twee stapjes werken. We gaan eerst eens kijken bij twee 98 Kelvin en één bar en dan gaan we eens kijken bij 500 Kelvin en tien bar.
1:31:21
Ik heb voor de volledigheid nog eens het startpunt uitgeschreven.
1:31:36
Dat is die gelijkheid bij evenwicht tussen de evenwicht constant en het reactie quotiënt.
1:31:40
Opnieuw, je kan een keuze maken op vlak van activiteit model.
1:31:46
Je kan ook werken met een partiële druk model bijvoorbeeld.
1:31:50
Je moet alleen zien dat je activiteiten waarin je de reactie quotiënt uitdrukt passen bij de
1:31:53
activiteiten die je berekent of die je gebruikt voor het berekenen van de evenwichts constanten.
1:31:58
Als we in een mol fractie model werken reactie quotiënt is die combinatie van activiteiten
1:32:03
waarbij elke activiteit een mol fractie wordt die altijd gebruikt wordt voor zo'n reactie.
1:32:09
Quotiënt producten in de teller telkens activiteit maal of tot een macht gegeven door een meters efficiënt reagentia in de noemer telkens
1:32:14
een activiteit tot een macht gegeven door de stok geometrische code efficiënt en dan bij evenwicht gelijk aan evenwicht constant.
1:32:22
In het mol fractie model is die evenwicht constant in een partiële druk model.
1:32:29
Het nulletje daar verwijst naar,
1:32:34
die kan je rechtstreeks als een exponentiële min delta geen nul op RT gaan bepalen en dan moet je eventueel corrigeren,
1:32:36
afhankelijk van de waarde van de druk en de waarde van die delta.
1:32:42
Nu ga je die moeten corrigeren om een X te berekenen.
1:32:45
De delta. Nu voor alle duidelijkheid is de sommatie komt van die substitutie van totaal doorrekenen.
1:32:49
Is de sommatie over alle producten van de documentaire significant minder sommatie over alle reagentia van een stoot meters efficiënt?
1:32:55
En daarnet die stikstof plus die zuurstof geeft twee mol.
1:33:05
En o, was die delta nu gelijk aan nul en je hebt twee als een soort gemiddelde hoeveelheid aan de product kant.
1:33:09
Één plus één is ook twee aan de reagens kant en dat geeft nul.
1:33:15
Hier heb je een voorbeeld waarbij je geen nul gaat krijgen,
1:33:19
dus tot je meters is één aan de ozon kant en is 3,5 aan de kant van de reagentia aan de zuurstof kant.
1:33:21
Misschien voordat we verder gaan. Wat gebeurt er met die hele set relaties?
1:33:33
Moest ik die toch met drie in plaats van drie en half en één als drie en twee schrijven?
1:33:41
Wat verandert er allemaal? Zou dezelfde evenwichts samenstelling moeten opleveren.
1:33:51
Hebben wij die de stootte met drie andere schrijven?
1:33:57
Ja, dat kan nooit een werkelijke evenwichts samenstelling kunnen veranderen.
1:34:01
Laat ons eens kijken hoe dat doorsijpelt in de relaties.
1:34:06
Als we die stof gemiddelde koffie zien te verdubbelen dan zie je dat hier elke NB maal twee gaat en elke nu gaat ook maal twee gaan.
1:34:10
Dus dat hele reactie quotiënt wordt ge kwadrateren. Een evenwicht constante.
1:34:19
De delta nu gaat ook twee keer zo groot worden, dus die term wordt ook kwadratisch in evenwicht constant.
1:34:25
Wat was daar de definitie van? KPI nul is exponentiële van min delta geen nul gedeeld door R2 en de delta geen nul.
1:34:33
Op zijn beurt is sommatie over alle producten storingsvrij meters efficiënt verwarmings vrije enthalpie van
1:34:44
dat product min sommatie over alle reagentia strooi meters efficiënt vorming vrije enthalpie van wat reagens.
1:34:53
Je ziet als ik al die stof meters efficiënter verdubbel dat je ook een twee keer zo grote delta geen nul gaat hebben
1:35:03
en dat is ook het linkerlid kwadrateren wordt en dus je mag zelf kiezen welke stoot je met drie je gebruikt.
1:35:10
Maar je moet natuurlijk consequent zijn en je moet je reactie consumenten en je
1:35:19
evenwicht constant en met dezelfde systematische structuur met drie gaan uitwerken.
1:35:24
En als je hier drie en twee schrijft dan moet je je delta geen nul gaan berekenen maar een twee en een drie.
1:35:29
Als je hier drie en half en één schrijft dan ga je delta geen nul maar half zo groot
1:35:36
zijn omdat je één gebruikt en drie en half gebruikt in plaats van twee en drie.
1:35:42
En dus de keuze van slots met drie is vrij blijft gewoon tot een kwadrateren of tot de vierkantswortel nemen links en rechts.
1:35:47
Maar je moet natuurlijk wel je reactie quotiënt en je evenwicht constanten consistent met dezelfde stof met drie gaan uitwerken.
1:35:55
Wat verkies je? Drie en twee of drie? Tweede en één? We werden met 23.
1:36:06
Goed. Wat krijgen we als reactie? Coachend? Wie? Wie werkt dat reactie coachend concreet uit?
1:36:20
Wie zegt wat ik moet neerschrijven? Als ik hier die breuk streep trek het gelijkheid teken.
1:36:25
En de evenwichts constanten. Wat gaat er in het rechterlid komen? Voila,
1:36:33
activiteiten van al uw producten tot een macht gelijk aan u terug in de teller
1:36:57
activiteiten van de reagentia tot een macht gelijk aan de gemiddelde groei in de noemer.
1:37:03
Je ziet, ik heb weer één vergelijking in in dit geval twee onbekenden.
1:37:07
Kan je oplossen door een massa balans in te voeren. Als je maar twee onbekenden hebt, kan je ook een beetje een kortere route nemen.
1:37:13
Mol fractie van AA is en AA gedeeld door A+ en B.
1:37:23
Mol fractie van B is en B gedeeld door en A+ en B.
1:37:31
Dus de som van de Mol fracties. Is altijd gelijk aan één.
1:37:38
Als je maar twee mol fracties hebt, dan kan je de ene altijd schrijven als één van de andere,
1:37:47
dus je mag daar nog altijd een massa balans aan toevoegen, maar hier is dat eigenlijk een beetje overbodig.
1:37:53
Het is pas als je met drie mol fracties werkt dat je al meer goed moet gaan nadenken hoe
1:38:00
je via dat soort relaties de ene mol fractie in functie van de andere kan gaan schrijven.
1:38:04
En dan wordt de massa balans heel handig om geen fouten te maken.
1:38:08
Als je maar twee mol fracties hebt is het een beetje overdreven om daar een ganse tabel voor te gaan neerschrijven.
1:38:12
We kunnen gewoon de Mol fractie, in dit geval de Mol fractie omschrijven als één min de mol fractie zuurstof of omgekeerd.
1:38:18
Ik ga er voor kiezen om de Mol fractie Ozon te laten staan.
1:38:25
En dan de Mol fractie. Zuurstof krijg je geschreven als één de Mol fractie Ozon en dan gaan we nemen tot de derde macht.
1:38:30
Hoe komen we aan die kicks? Dat is de volgende stap.
1:38:45
Tot nu toe of voorbij 500 Kelvin of tien en tien bar werken of bij twee 98 Kelvin en één bar.
1:39:02
De relatie is dezelfde. De specifieke omstandigheden waarbij we werken gaan zich vertalen in het berekenen, in de manier waarop we X gaan berekenen.
1:39:09
Stellen. We starten bij twee 98 Kelvin één bar. Hoe kunnen we X gaan bepalen?
1:39:20
Wat is de algemene relatie? Wat is de algemene relatie en hoe vertaalt zich die concreet als we met die omstandigheden werken?
1:39:26
De algemene staat op bord is de KP nul en dan P totaal over P tot de macht minder dan nu.
1:39:41
Als we het over een één bar gebruiken, waar gaat die X aan gelijk zijn?
1:39:49
Wat gebeurt er met deze term? Als we onder die omstandigheden werken.
1:39:59
We gaan gewoon KPI nul zijn. Happy Totaal is op dat moment gelijk aan nul één, dus je moet het dan nu al niet meer berekenen.
1:40:20
Dat gaat altijd één opleveren. Dus op dat moment is je X gelijk aan je KP nul.
1:40:27
BTW 98 Kelvin.
1:40:32
Dat betekent dat je Delta geen nul gewoon uit de linkse vrije enthalpie uit de tabel kan plukken, want dit is een tabel met twee 98 Kelvin.
1:40:33
Je delta geen nul. Is in dit geval hoeveel keer de ingrediënten van ozon.
1:40:43
Voor de gekozen fotogrammetrie moeten we. Hoe gaan we twee keer.
1:41:05
De vorming is vrij enthalpie. Hebben van ozon -3 keer de vorming vrije enthalpie van CO2.
1:41:10
En dan kan je je getallen beginnen aflezen. Voor ozon heb je daar 163,2 maal twee.
1:41:20
Dat gaat je honderd 300. 26 als ik dat snel uitreken kom maar vier opleveren kilo per mol.
1:41:27
Die zuurstof is de zuivere enkelvoudige stof. Het mag dus niet verbazen.
1:41:39
De vorming is vrij tien nul is. Eigenlijk is dit gewoon de vorming reactie van Ozon,
1:41:43
maar dan twee keer genomen en dus de delta geen nul is twee keer de vorming vrije enthalpie van ozon 326,4 kilo zoelen per mol.
1:41:49
Wie berekent er de evenwicht constanten daarmee? Wacht even.
1:42:02
163. Maal twee is volgens mij 326,4.
1:42:11
Ja oké, reken eens uit. Wat heb je van enig constante?
1:42:22
Ik wil een getal horen en ik wil daar geen half uur op moeten wachten.
1:42:29
Hoe kom je dan aan? Drie.
1:42:49
Ja, dat hebben we al gedaan. Het is de vorming van vrije wil, die van ik wilde reactie vrijheid.
1:42:52
Ik wilde even de constanten kennen.
1:42:57
Dus de exponentiële van Delta geeft nul op de delta er geen nul is twee keer de vorming van oceanen en dus ik wil nu de.
1:43:01
Ik wil een getal. Je doet toch wel dezelfde fout als vorige week gemaakt.
1:43:17
En natuurlijk, ja, je hebt iets van 326.000 en je gaat dat delen door 8,31 maal 300.
1:43:54
Ja, dat delen door ongeveer 3000 is wel een zeer kleine toestand.
1:44:01
Hoeveel kom je uit? Het was wel een heel klein.
1:44:05
Maar het kan. Ik heb het nu zelf uitgerekend, om het zo maar te zeggen.
1:44:17
Dat is te klein. Laat ons. Ik zal er zelf even aan beginnen. Dus we hadden 300 26.400.
1:44:21
Ola. Gedeeld door 8.314 gedeeld door 200 98 is 131.
1:44:41
Dus je gaat gelijk hebben en dan krijgen wij met een minteken erbij hier.
1:44:56
Voila, en dan moet ik de exponentiële nemen. En ik denk dat ik dan.
1:45:02
Second hier. Voila.
1:45:09
Dus we kregen iets bijzonder klein. Zes. Tien tot de -58 te.
1:45:12
Kan. Waar ik wil dat je op let.
1:45:23
Heel vaak gaan die vergelijkingen je geen lineaire of kwadratische vergelijking opleveren.
1:45:47
Ik kreeg hier een mooie fractie in het kwadraat.
1:45:55
Ik krijg een mol fractie tot de derde, dus dat is iets dat je niet gaat kunnen omzetten naar een lineaire kwadratische vergelijking.
1:45:58
Dus je gaat nooit met de hand exact een oplossing kunnen vinden.
1:46:06
Wat je moet proberen begrijpen is. We hebben hier een heel klein getal aan de rechterkant.
1:46:11
De fracties die we zoeken, die liggen tussen nul en één.
1:46:20
Als de Mol fractie gelijk is aan nul dan krijg je nul gedeeld door één -0.
1:46:26
Dan heb je nul aan de linkerkant. Als de fractie van Ozon één wordt, dan krijg je hier één gedeeld door één -1.
1:46:32
Dat gaat naar oneindig aan de linkerkant. Het rechterlid in dit geval.
1:46:43
Is zeer klein. Om een zeer klein rechterlid te hebben gaan we hier een teller moeten hebben die zeer klein is en een noemer die dan gelijk is aan één.
1:46:50
We gaan met andere woorden oplossingen gaan vinden waarbij die mol fractie van O praktisch gelijk is aan nul.
1:47:01
Omgekeerd zou je kunnen hebben dat je hier een heel groot getal uitkomt.
1:47:10
Dat gaat alleen maar lukken op het moment dat je mol fractie van ozon in de richting van één gaat,
1:47:14
want je hebt dan één gedeeld door iets dat de lucht in. Dan kan één net zo groot worden als je maar wilt.
1:47:20
Hoe dichter je bij één komt, hoe groter dat gaat worden. Dus je moet die relaties proberen te benaderen vanuit wat zijn de twee limiet situaties?
1:47:26
Want je gaat altijd iets hebben van een mol fractie of een vordering graad en één vordering graad.
1:47:35
Als die vordering schade klein zijn dan kan je één min altijd gaan benaderen als de vorige factor.
1:47:43
Als die eerste term als die één en dus de mol fractie bij van ozon.
1:47:50
Bij evenwicht ga je in dit geval berekenen uit kwadraat van die mol fractie gedeeld door één.
1:47:55
In principe zou daar staan waarschijnlijk één min iets dat gelijk is aan 0,0000000.
1:48:02
En dan ga ik nog waarschijnlijk twintig nullen verder moeten gaan voordat je een één hebt,
1:48:09
maar dat is in wezen gelijk aan één en dus benaderend vind je?
1:48:13
En dat is in dit geval een zeer goede benadering. Dat die mol fractie van Ozon de vierkantswortel gaat zijn uit zestien tot de -58.
1:48:17
Wat betekent dat je iets gaat hebben van de jaren. Wortel zes.
1:48:33
Tien -29. Hoe kunnen we daar een interpretatie aan geven?
1:48:39
Een ander is zelfs eigenlijk geen ozon,
1:48:49
want in één mol zit je met tien tot de 23e deeltjes en je hebt 1 miljoen mol nodig om één molecule tegen te komen.
1:48:51
Dus dan kan je eigenlijk zeggen in alle praktische toepassingen ga je bij kamertemperatuur in evenwicht.
1:49:01
Met zuurstof gaat er nooit ozon gevormd worden.
1:49:06
Als we naar andere temperaturen en andere drukken gaan, dan gaat die constante wijzigen en kan die samenstelling veranderen.
1:49:12
Dus laat ons eens kijken wat er gebeurt als we de temperatuur opdrijven en als we de druk opdrijven.
1:49:23
Dus opnieuw opletten als je even iets constante berekent. Invullen in joule per mol en niet kilo vullen per mol.
1:49:49
Als je een constante hebt die ongeveer één is, moet er altijd een reactie zijn.
1:49:55
Ik zit waarschijnlijk een factor duizend verkeerd met mijn delta nul.
1:50:00
Dus X. Is.
1:50:12
Mol fractie drie kwadraat gedeeld door één min. Mol fractie drie tot de derde macht is de KP nul Mol p.
1:50:17
Totaal op nul tot de min delta nu.
1:50:29
We gaan niet bij tien bar werken. We gaan bij 0,1 bar werken om toch een beetje de goede richting uit te gaan.
1:50:36
Bepalen we nu de KPI nul. Als ik terugkeer naar.
1:50:58
Er is een tabel met twee 98 kelvin. Hoe kan je de data gebruiken?
1:51:16
Om toch een delta geen nul bij 500 Kelvin te bepalen. We hebben al een paar keer.
1:51:22
Gehad. Wie zegt het nog eens uit?
1:51:28
Je haalt de Delta H0 en de delta snel uit de tabel en dan gebruik je die waarden om bv over een delta geen nul te gaan bepalen.
1:51:33
Dus voor Delta H0. Ga je twee keer de vormings enthalpie hebben van ozon.
1:51:40
Want opnieuw zuurstof is je referentie toestand, dus je krijgt daar een Delta H0 van 285,4 kilo joule voor mol en je Delta S nul.
1:51:51
Gaat het verschil zijn tussen die absolute entropie is dan twee keer.
1:52:07
De entropie van ozon is twee maal 238,9.
1:52:16
Drie keer die van zuurstof is, terwijl zijn vijf.
1:52:26
kilo. Jules per mol, Jules per Kelvin per mol. En Delta G nul ga je dan berekenen als delta nul -500 maal de delta is nul.
1:52:32
O ja, dit is beter daar. Dat is de keurige notatie.
1:53:06
Sorry. Wie heeft er waarde voor het geld dat geen nul.
1:53:12
Op. Ik heb 354.300.
1:54:58
Jules per Mol. Jij ook?
1:55:48
Ja of nee. Oh ja, dat kan zijn dat ik ook misschien ergens.
1:55:54
Dat is toch wel heel erg in de buurt. Je hebt het ook oké. Dus wat krijg je van KPN nul?
1:56:06
Moet je delen door R en in dit geval ga je delen door 500 Kelvin in plaats van twee 98 Kelvin.
1:56:13
Dus we gaan dit getal delen door. Er.
1:56:21
En we gaan delen door 500.
1:56:28
We gaan naar een minteken en we doen de exponentiële.
1:56:34
En nu heb ik 9,6. Tien -38 is nog altijd klein, maar ik wil je aandacht vestigen op.
1:56:41
Iets dat belangrijk is. Je ziet dat de delta geen nul.
1:56:52
Groter is. Dan Degene die we berekenden werd twee 98 Kelvin en daar hadden we iets van 326 en zoveel.
1:56:58
Maar de enige constante is kleiner. Onthou de constante wordt bepaald door de verhouding tussen delta geen nul en T.
1:57:07
In dit geval geldt dat genot stijgt.
1:57:16
Ja, maar het stijgt nog meer waardoor dat de exponent omlaag gaat en je in dit geval een minder kleine, een iets grotere constante berekent.
1:57:19
Dus even flink hangt af van groei op tv en niet zozeer van groei alleen.
1:57:30
Dat heeft zijn impact. Dat gaan we zien op de temperatuurs afhankelijkheid van zo'n evenwichts relatie.
1:57:36
Hier heb je een voorbeeld. Als ik naar hogere temperaturen ga Delta geen nul wordt nog meer positief.
1:57:43
Maar dat betekent niet dat de constante nog kleiner wordt.
1:57:50
Omdat we delen door een nog grotere temperatuur gaat de constante eigenlijk groter worden.
1:57:53
Ze blijft klein, maar ze is een pak groter, ongeveer twintig grootteorde groter dan wat we daarnet hadden.
1:57:59
Hoe verwerken we die totaal? Druk. 0,1 bar.
1:58:11
Daar moeten we hier in mee nemen. Die totaal is nu 0,1.
1:58:15
Het is altijd pi totaal gedeeld door standaard. Truc is 0,1.
1:58:21
De delta nu zoals we daar geschreven hebben is -1,
1:58:25
dus dat gaat ons nog een factor tien opleveren voor de evenwicht Constante
1:58:28
0,1 tot het eerste is nog eens een factor tien voor de evenwichts constante.
1:58:33
Dus de X is de KP nul maal tien.
1:58:39
Dus de X in dit geval gaat 9,610 -37 gaan bedragen en dat is opnieuw gelijk aan.
1:58:42
Je ziet dat je nog altijd in de limiet kan werken om zo'n klein getal te hebben gaat die fractie van Ozon eigenlijk nul gaan zijn.
1:58:52
En in de fractie van Ozon gaan we nog altijd aan één kunnen gelijkstellen, dus dat gaat ongeveer.
1:58:59
Maar die ongeveer is praktisch gelijk aan de fractie van Ozon in het kwadraat gaan zijn.
1:59:05
Je pakt links en rechts de vierkantswortel en dus in dit geval krijg je een fractie van ozon die tien -18 bedraagt.
1:59:12
Maar dat betekent dat je op één mol van die moleculen al honderdduizend ozon moleculen gaat hebben.
1:59:20
Dus je voelt aan dat zo'n chemische mislukking dat die heel sterk kan gaan wijzigen als de temperatuur verandert.
1:59:25
Natuurlijk blijft de Mol fractie van Ozon zeer laag, maar is enorm gestegen vergeleken bijvoorbeeld bij 98 Kelvin.
1:59:32
Ik vermenigvuldig met tien. We hebben een totaal is 0,1, maar tot de macht minder dan nu.
1:59:46
Delta nu is. -1.
1:59:54
Wacht efkes. Je hebt gelijk. Je moet.
1:59:57
Ik weet totaal. Is 0,1.
2:00:00
Tot de macht. Minder dan nu. Minder dan nu is -1.
2:00:06
Waar is Tina? Oh. Dus we gaan uit van tien -38 naar tien -37 nog eens bijna maal tien dus is tien -36 vierkantswortel tien -18.
2:00:10
Dus we leren een aantal dingen. Dat is ten eerste hoe we die berekeningen moeten doorvoeren.
2:00:23
Doe dat stapsgewijs. Eerst KP nul. Gebruik de juiste stokjes met drieën.
2:00:28
Gebruik de juiste temperaturen. Ga van KP nul naar X en dan kijk je of je zo reactie quotiënt kan vereenvoudigen.
2:00:32
Heb je argumenten om te zeggen die mol fracties die we zoeken die zijn hetzij praktisch nul, hetzij praktisch één, Dan ga je termen kunnen wegstrepen.
2:00:40
In dit geval Mol fracties zijn praktisch nul. Deze gaat over overblijven.
2:00:51
Die gaat eigenlijk één worden en je vindt vrij gemakkelijk oplossingen zonder dat je die ganse vergelijking moet gaan oplossen.
2:00:56
Als je er echt niet uitgeraakt kan je ook altijd met je rekenmachine proberen dat op te lossen.
2:01:02
Dat gaat doorgaans wel lukken. Het is altijd goed om zelf die limieten in de gaten te houden,
2:01:07
want het is nogal snel dat je een typfout gemaakt hebt in het ingeven van data in je rekenmachine.
2:01:12
Net zoals je factor duizend die je mist in het berekenen van evenwicht constanten.
2:01:18
Dus zelfs als je gewoon via je rekenmachine probeert de wortels van zo'n vergelijking te vinden,
2:01:22
hou in de gaten of die wortels die je eruit haalt effectief wel zin hebben en of je geen fouten gemaakt hebt bij het invullen van die getallen.
2:01:26
Het is daarom dat dat inschatten van oplossingen wel leerrijk is om te doen.
2:01:33
Het tweede dat we leren is dat die evenwicht constant en daardoor ook evenwichts
2:01:37
liggingen sterk kunnen afhangen van druk en sterk kunnen afhangen van temperatuur.
2:01:41
Dat is waar ik nog even dieper op wil ingaan, omdat je daar een aantal systematische verbanden naar voor kunt schuiven.
2:01:46
De manier waarop we activiteiten van componenten in oplossing bepalen.
2:01:59
Ik denk dat we dat volgende les gaan zien, tenzij ik seffens nog tijd genoeg heb.
2:02:03
Even geduld. Daar zijn we, dus dat gaan we even achterwege laten.
2:02:07
Misschien dat ik nog even daarop terugkom. Maar ik wil naar hier gaan.
2:02:14
Dat is invloed van temperatuur op een chemische evenwichts ligging wat we hier al kort hebben gezien.
2:02:20
Dus ze zijn gewend ondertussen om te denken in termen van evenwicht betekent reactie
2:02:39
quotiënt in een bepaald activiteit model moet gelijk zijn aan de evenwicht constante die bij
2:02:45
dat activiteit model hoort en we kunnen activiteiten model kiezen waarmee we het evenwicht
2:02:52
constant en rechtstreeks uit de thermodynamische tabellen kunnen halen en dan schrijven,
2:02:56
als dat het geval is, een nuttige boven die evenwicht constanten.
2:03:01
Ik denk dat je ook begrepen hebt dat. Hele kleine is dit constante doorgaans betekent nauwelijks product.
2:03:05
Zeer kleine activiteiten voor het product vergeleken met de activiteit voor de reagentia.
2:03:14
Dat is de situatie die we hier hadden. Hele kleine constante levert nauwelijks iets op op vlak van ozon.
2:03:20
Omgekeerd hele grote evenwicht Constante betekent doorgaans hoge activiteiten voor je product en kleine activiteiten voor je reagentia.
2:03:27
Je moet een beetje daarmee voorzichtig zijn, omdat grote of kleine evenwicht constanten afhangen van de definitie van je standaard toestanden.
2:03:37
Bijvoorbeeld door de druk te veranderen gaat de KPI nul die je hier hebt.
2:03:46
Die gaat ongewijzigd blijven. Maar een evenwicht kan wel verschuiven als we de druk heel erg verhogen of heel erg verlagen.
2:03:52
Dus de waarde van een evenwicht constante is niet noodzakelijk direct een vertaling
2:04:01
van de ligging van een evenwicht waar je wel vrij veilig kan aannemen of kan geven.
2:04:05
Kan van uitgaan is als een evenwicht constant groter wordt.
2:04:10
Dus ik heb het niet zozeer over de absolute waarde, maar over de verandering van een evenwicht constante.
2:04:15
Als in evenwicht constante groter wordt, dan ga je hogere activiteit hebben van een product.
2:04:20
Typisch dus meer product en lagere activiteiten van een reagens.
2:04:25
Als een evenwicht constant kleiner wordt, gaat het omgekeerd zijn.
2:04:30
En dus Waar we naar gaan kijken is niet zozeer de absolute waarde van de evenwicht constante.
2:04:34
Waar we naar gaan kijken is hoe de evenwicht constante verandert.
2:04:39
Bijvoorbeeld als we de temperatuur wijzigen. En dan kan ik inschatten als ik die evenwicht konstante koppel aan de delta geen nul van een reactie.
2:04:43
En daar kom ik terug op het feit dat de constante de temperatuur afhankelijkheid zowel in de teller als in de noemer zit.
2:04:56
In plaats van direct een afgeleide van naar de temperatuur te doen, gaan we eerst gaan kijken naar de logaritme van de evenwicht constanten.
2:05:06
En eigenlijk de temperatuurs afgeleide van die logaritme proberen in te schatten.
2:05:17
Je weet een logaritme is een monotone functie van zijn argument.
2:05:23
Daarmee bedoel ik stijgt. Logaritmisch, stijgt als het argument stijgt.
2:05:27
Dus als Ellen omhoog gaat, maar de temperatuur gaat ook omhoog gaan met de temperatuur.
2:05:32
Dus of we dan temperatuurs afhankelijkheid van de evenwichts constanten kijken
2:05:38
of de temperatuurs afhankelijkheid van de ellen van de eeuwig constante.
2:05:41
Dat is eigenlijk hetzelfde als de zelfde informatie die we daaruit halen. Waarom is het interessant om naar die ellende te gaan kijken?
2:05:45
Wel, als je de temperatuurs afgeleide neemt en ik gebruik een partiële afgeleide omdat we werken bij vaste druk bijvoorbeeld,
2:05:52
dan krijg je hier een temperatuurs afgeleide ook.
2:06:00
En dat is een relatie die we zijn tegengekomen en waarvan ik gevraagd heb.
2:06:07
Bekijk die zelf een keer en zie dat je die kan afleiden.
2:06:14
Ik ben niet zeker of je dat gedaan hebt, maar wat we hier hebben is temperatuurs afgeleide van G gedeeld door T.
2:06:19
Waaraan is dat altijd gelijk een temperatuurs afgeleide van een g gedeeld door T.
2:06:25
WiFi afleiding proberen maken. En dan is er nog een tweede voorwaarde voor het resultaat te houden.
2:06:38
Dus ik ga ervan uit dat jullie allemaal die aflevering gedaan hebben, maar dat je het resultaat vergeten hebt of de temperatuur zakt.
2:06:45
Verder van geopteerd als de Gibbs Helmholtz relatie. Ik ga het nog eens zeggen.
2:06:54
Probeer die af te leiden en probeer ze niet alleen af te leiden, maar probeer ook het resultaat te onthouden en geef op die af.
2:06:59
Lennart geeft een min ha gedeeld door t kwadraat.
2:07:05
Dus als ik hier een dictee heb van G gedeeld door twee dan kan ik dat schrijven als delta nul van de reactie gedeeld door RTE kwadraat.
2:07:12
Wat je ziet is. Het teken van die afgesneden.
2:07:25
Stijgen van de constante bij stijgende temperatuur gaat afhangen van het teken van de delta.
2:07:32
Er is altijd positief. Kwadraat is altijd positief.
2:07:40
Het enige wat ervoor kan zorgen dat het rechterlid hetzij positief, hetzij negatief is, is de delta.
2:07:43
Voor een reactie. Waar de delta haar positief is. Dat is een reactie die warmte opneemt.
2:07:51
Een endo terme reactie hebben dat de constanten groter wordt als de temperatuur omlaag gaat.
2:07:56
Voor reacties die warmte afgeven en negatieve delta externe reacties gaat de constante kleiner worden als
2:08:04
de temperatuur omhoog gaat en dus voor andere termen reacties zoals bijvoorbeeld die vorming van ozon.
2:08:12
De delta is plus 285 kilo joule per mol.
2:08:19
Temperatuur opdrijven betekent richting.
2:08:23
Uit het evenwicht verschuift de evenwicht, constante wordt groter, Het evenwicht verschuift in de richting van de producten.
2:08:26
Je hebt hier ook gezien de evenwichts mol fractie van ozon gaat omhoog, blijft klein maar gaat omhoog als we van 290 naar 500 Kelvin gaan.
2:08:32
Voor extreme reacties geldt het omgekeerde.
2:08:43
Als je daar de temperatuur opdrijft, dan gaat de evenwichtige uit het evenwicht in de richting van de reagentia gaan verschuiven.
2:08:45
Algemeen kan je dat samenvatten als temperatuur Opdrijven doet een evenwicht verschuiven in interne
2:08:52
zin in de richting die warmte opneemt voor een andere term reactie is dat de voorwaartse richting,
2:08:58
voor een externe reactie is dat de achterwaartse richting. Dan heb je misschien geleerd als het principe van licht atelier,
2:09:04
maar ik vind het beter om dat te onthouden aan de hand van de tekens van die afgeleiden,
2:09:12
omdat het principe van Locatelli werkt als je het juist formuleert en niet werkt als je het verkeerd formuleert.
2:09:16
Terwijl die constanten, die kan je maar op één manier gaan schrijven en daar kan je geen vergissingen in maken.
2:09:22
We hebben hier ook de druk veranderd en gekeken hoe dat dat in evenwicht beïnvloedt.
2:09:37
Het is een beetje subtieler. KP nul.
2:09:47
Gedefinieerd als exponentiële van min delta geen nul.
2:10:26
Gedeeld door RTE. Opnieuw zou ik naar de ellende kunnen gaan.
2:10:32
Wat gebeurt er met het linker lid? En wat gebeurt er met het rechterlid?
2:10:52
Dus. Als ik al naar de druk.
2:10:58
Rechterlid is wel dat geen nul is, dus de reactie vrije enthalpie bij een druk van één bar gedeeld door RT.
2:11:12
En we gaan eens kijken wat er gebeurt als we de reactie doorvoeren.
2:11:20
Niet bij één bar, maar bij 0,1 bar zoals we daarnet gedaan hebben.
2:11:24
Wat gebeurt er dan met die KP nul? Het is een beetje een strikvraag.
2:11:29
Helemaal niets. Die Capitole is gedefinieerd bij één bar, zoals we daarnet ook gedaan hebben als we werken bij 0,1 bar.
2:11:49
We hebben de KPI nul berekend op basis van tabel.
2:11:57
Hier de waarden. De tabel staat er niet meer, maar we zijn naar die tabel gaan kijken om de KPI nul te berekenen.
2:12:01
De druk afhankelijkheid hebben we pas later verwerkt als we van de KP nul naar de X gegaan zijn.
2:12:07
Dus hier heb je een voorbeeld waarbij je moet opletten met de waarde van een evenwicht.
2:12:14
Constante KP nul is druk onafhankelijk.
2:12:18
Maar bijvoorbeeld een evenwicht zoals de ozon vorming kan wel degelijk gaan verschuiven als je de druk verandert.
2:12:22
Dus uit. De waarde van een leven is constant en zelf moet je altijd een beetje voorzichtig zijn met conclusies te trekken.
2:12:29
Op vlak van ligging van een evenwicht, bijvoorbeeld een KP nul, ga je altijd dezelfde constanten opleveren.
2:12:34
Maar als je dat bericht bekijkt bij tien -6 bar of bij tien tot de zesde bar kan de afwisseling heel erg verschillen.
2:12:41
Om die druk afhankelijkheid te laten uitkomen is het beter in dit geval naar de X te gaan kijken.
2:12:48
X hadden we gedefinieerd als KP nul en dan kwam daar die P totaal gedeeld door P nul tot de min.
2:12:56
Dat dan nu bij kijken ook. En dus als je wilt zien hoe een evenwicht verschuift in functie van de druk is de X een veel betere insteek dan de KP nul.
2:13:04
De KP nul als waarde is blind is onafhankelijk van de druk.
2:13:14
Als ik hier naar Ellen X ga. Opnieuw is dat interessanter dan de X zelf afleiden.
2:13:19
Krijg je Ellen van de KPN nul en kreeg je min delta nu Ellen van P.
2:13:25
Totaal gedeeld door standaard druk. En dus als we afleiden naar de druk.
2:13:31
Eerste term hebben we al gehad. Die is druk, onafhankelijk.
2:13:42
Dus als we afglijden naar de druk verdwijnt die.
2:13:45
En dan hou je over minder dan nu en dan ellende van het totaal op nul geeft je P nul gedeeld door P totaal.
2:13:48
Als afgeleide. Wat doe ik daar allemaal?
2:13:59
Is één op nul. Ellen is gewoon één op één totaal hè?
2:14:07
Ja, voilà, dat is het. De Ellen van één op nul valt weg, is een constante.
2:14:14
Is dat dan nu gedeeld door P Totaal? Je ziet, het teken gaat opnieuw afhangen.
2:14:19
Niet van die P totaal. Die is altijd positief, maar gaat afhangen van het teken van de delta.
2:14:25
Nu, dat speelt dezelfde rol als exo term of andere term karakter.
2:14:30
Als Delta nu. Kleiner is dan nul of telt dan nu?
2:14:36
Is groter dan nul voorstelt dan nu kleiner dan nul, dan gaat een constante constant te gaan stijgen.
2:14:41
Als de druk verhoogt wordt delta nu groter dan nul, gaat het evenwicht constant te gaan dalen als de druk verhoogd.
2:14:47
Is samengevat in die tabel. En dat dan nu positief.
2:15:02
Geeft je een DLM op DP kleiner dan nul?
2:15:06
Zware ja wat ik daar ook heb geen contradicties.
2:15:10
Delta nucleaire nul geeft je een die alleen op DP groter dan nul indien in dit geval als de druk omhoog gaat.
2:15:14
Evenwicht gaat in de richting van de reagentia. De X die daalt best bij stijgende druk.
2:15:22
B omlaag even gaan in de richting van de producten en omgekeerd.
2:15:28
Opnieuw kan je het samenvatten als je de druk verhoogt, gaat en gaat en gaat het evenwicht verschuiven in de richting van minder mol stof.
2:15:31
Als de druk verlaagt, gaat het evenwicht gaan verschuiven in de richting van meer mol stof.
2:15:39
Heeft zijn toepassing gehad. Dat is misschien waar ik de les mee ga proberen beëindigen.
2:15:48
De figuur die daar was heb ik niet. Wacht even, het is hier vastgelopen.
2:16:03
Ik ga er een tweede voorbeeld bijhalen. Even geduld.
2:16:55
Ja. Oké, ik wil er een tweede voorbeeld bijhalen.
2:17:44
Toch wel van historisch belang geweest is. Laat ons eerst eens proberen die inzichten zelf uit te schrijven.
2:18:44
Het gaat over de vorming van ammoniak. Uit stikstof en waterstof.
2:18:50
Als je dat goed leest is dat opnieuw de voedings reactie van ammoniak.
2:19:02
We vormen ammoniak uit stabiele en enkelvoudige stoffen die elk dienen als referentie toestand, hetzij voor stikstof,
2:19:06
hetzij voor waterstof, is historisch gezien een belangrijke reactie geweest, Eigenlijk omwille van twee zaken.
2:19:14
Ammoniak is een basisproduct van de chemische nijverheid, onder andere voor de productie van allerhande kunstmeststoffen, Nitraten bijvoorbeeld.
2:19:25
En in de. Negentiende eeuw.
2:19:36
Ja, heb je. Een sterke toename gehad van de bevolking is een bevolkingsexplosie die eigenlijk gevolgd is op de industriële revolutie.
2:19:40
En op het moment was er wel een probleem op vlak van voedselproductie en op vlak van productiviteit van de landbouwgrond eigenlijk,
2:19:49
en men moest daar nitraten voor hebben die konden dienen als kunstmest,
2:19:57
eigenlijk om per hectare meer voedsel te produceren en de enige bron van nitraten was.
2:20:02
Vogel mest En daar kom je waar een aantal eilanden in de Stille Zuidzee,
2:20:10
waar van die enorme broed kolonies zitten en waar tientallen meters vogel mest vogels stront verzameld liggen.
2:20:15
En die werden ook opgekocht. Dat was een soort business om die eilanden te kopen en daar die die mest te gaan, te gaan verzamelen.
2:20:22
Te gebruiken als basis, als startpunt van een hele nitraat industrie.
2:20:31
Dat is natuurlijk niet zo praktisch en waar men naar streven is om eigenlijk dezelfde chemie te kunnen opstarten,
2:20:37
maar dan op basis van ammoniak en wat ook en een stikstof houdende component is die je in het beste geval kan synthetiseren uit stikstof en waterstof.
2:20:44
Stikstof dat je gewoon uit de lucht kan halen en dan kan je die ganse nitraat
2:20:53
verwerkende nijverheid opstarten zonder dat je last moet hebben van van die.
2:20:57
Die guano heette dat of heet dat nog altijd die je kan gaan gaan opscheppen in de in de Stille Zuidzee?
2:21:02
Waarom is dat niet zo'n eenvoudige reactie? Als je weet wat stikstof.
2:21:10
Wat voor elektronische structuur stikstof heeft.
2:21:18
Waarom is dat niet zo'n eenvoudige reactie? Een mengsel van stikstof en waterstof.
2:21:21
Dat kan je gewoon kopen. Bij Air Liquide kan je flessen kopen waar een mengsel van stikstof en waterstof in zit en dat mengsel blijft bestaan,
2:21:32
Dus niet dat dat omgezet wordt na verloop van tijd in een ammoniak.
2:21:41
Daar gebeurt eigenlijk niets mee. Waarom is dat zo?
2:21:44
Waarom is die omzetting? Want we gaan zien dat die ammoniak vorming thermodynamisch voordelig is.
2:21:48
Waarom gebeurt die niet zo gemakkelijk? Is het stabiel en kan je dan linken met de elektronische structuur van stikstof.
2:21:55
Was die driedubbele binding. Denk aan de voeding van die moleculaire orbitalen.
2:22:07
Die drie bindende orbitalen zijn gevuld.
2:22:10
Drie anti binnen de orbitalen zijn leeg, wat je kan vertalen in die drie dubbele binding in de logische structuur om ammoniak te vormen.
2:22:12
Er staat hier een voor factor. Een half is ook maar één atoom stikstof.
2:22:20
Ga je op een of andere manier die stikstof stikstof binding moeten verbreken?
2:22:24
Wat ervoor zorgde die reactie een behoorlijk grote activering Synergy heeft.
2:22:29
Voordat we naar de Kinetic gaan kijken wil ik eerst eens de thermodynamica en de evenwichts relaties voor de reactie gaan neerschrijven.
2:22:33
Opnieuw kunnen we die gaan analyseren in termen van evenwicht. Constante moet gelijk zijn aan reactie quotiënt.
2:22:41
Wat wordt het reactie quotiënt in dit geval? Wie goed naar de slides kijkt, ziet daar een suggestie.
2:22:49
Maar er staan ook foutjes op de slides, zie ik. Dus misschien beter om het op bord te zetten.
2:23:01
Hier heb je weer een voorbeeld. Geeft een vergelijking in drie mol fracties.
2:23:19
Dus het trucje de Mol fractie van de ene plus de mol fractie van de andere geeft één is
2:23:24
onvoldoende om dat te substitueren en dan één vergelijking in één onbekende van te maken.
2:23:29
Je kan proberen via de originele verbanden op te stellen.
2:23:34
Hier is het aangewezen om direct over te gaan op zo'n massa balans.
2:23:39
Laat ons dat nog eens bekijken. We hebben stikstof, we hebben waterstof, er is ammoniak, er is een totaal aantal mol dat we moeten meepakken.
2:23:43
Ik heb een vordering graad. In het begin is die nul. De vraag is wat er gebeurt als de vordering is graad alfa en nul bedraagt en
2:23:54
dan moet je een bepaalde begin concentratie of begin aantal mol voorstellen.
2:24:02
Dat zou bijvoorbeeld het metrische mengsel kunnen zijn. Daar hebben we 2.00u mol stikstof.
2:24:07
We hebben dan 3,5 en nul mol waterstof. En ik heb geen ammoniak, want in totaal twee en nul mol in het systeem geeft.
2:24:13
Je kan daar andere samenstellingen gebruiken, maar andere begin samenstellingen gaan andere evenwichts samenstellingen opleveren natuurlijk.
2:24:22
Hier start ik van iets dat in een opgave een stokje of metrisch mengsel zou genoemd worden.
2:24:29
Dat is dat de aantallen gelijk zijn aan dus of evenredig zijn met de structuur efficiënt in de reactievergelijking.
2:24:33
Als de omzetting graad alpha en nul bedraagt.
2:24:41
Wat verandert er aan de hoeveelheid stikstof? Op twee.
2:24:45
Voor de koffies die ik gekozen heb gaat dat over en nul gedeeld door twee opleveren.
2:25:03
Je gaat ook waterstof zien verloren gaan met een bedrag -3, tweede alfa en nul en er gaat ammoniak gevormd worden met een bedrag alfa en nul.
2:25:08
Je ziet in dit geval anders dan bij dat stikstof,
2:25:19
zuurstof en nul dat het totaal aantal mol ook gaat veranderen en verdwijnt twee alfa en nul mol en er komen maar één keer alfa en nul mol terug.
2:25:21
Dus je bent alfa en nul mol kwijt als de reactie alfa en nul keer is doorgegaan.
2:25:31
Er verdwijnt meer aan de linkerkant dan dat er bijkomt aan de rechterkant.
2:25:37
Totaal heb je 2.00u.
2:25:45
In mijn Alpha. Je hebt 3/2 en nul.
2:25:50
Één mijn alfa en je hebt alfa en nul keer.
2:25:56
Ammoniak. Totaal aantal mol is twee min alfa.
2:26:01
Nul. Daaruit kan je dus dat zijn de totale waarde van en daaruit kan je de mol fracties gaan berekenen.
2:26:06
Dat is die één nul gedeeld door nul. Dat valt weg.
2:26:13
Je hebt dan één min over gedeeld door twee keer twee min alfa voor ammoniak.
2:26:17
Je hebt een drie. Één min alpha gedeeld door twee keer twee min.
2:26:25
Alva voor waterstof en dan heb je een Alfa.
2:26:31
Gedeeld door twee mijn alfa voor het ammoniak.
2:26:37
En daar heb je die drie mol fracties in één enkele vordering geschrapt die ik kan gebruiken om die evenwichts relatie vorm te geven.
2:26:42
En dus de X. Neem me voorlopig over gaan we schrijven als de Mol fractie van ammoniak als alfa.
2:26:56
Gedeeld door twee in Alpha. Die van stikstof is één min alfa gedeeld door twee keer twee in alfa en dat is een vierkantswortel en dan
2:27:06
krijgen we drie keer één min alfa gedeeld door twee maal twee in alfa en dat is tot de macht drie halven.
2:27:19
Dan gaan we een beetje kunnen opkuisen. We houden een alfa daarboven.
2:27:31
Ik heb hier een twee min alfa. Ik deel door een wortel van twee mijn alfa.
2:27:36
Hier deel ik ook nog eens door een wortel van twee, mijn alfa dus dat gaat 2.01u mijn Alfa overbleven die naar de teller komt,
2:27:40
dus dat geeft me een twee mijn Alfa in de teller en dan heb je één mijn alfa, de wortel en één mijn alfa drie halve.
2:27:49
Geef mijn één in alfa kwadraat in de noemer en dan de factoren die zijn overgebleven is hier een drie tot de drie halven.
2:27:57
En heb ik hier een wortel twee en een twee tot de drie halve.
2:28:12
Dat is samen een twee kwadraat die ook boven terecht komt, dus dat gaat een vier zijn die daar staat.
2:28:16
Als ik mij alles goed onder controle gehouden heb is dat de uitdrukking voor de evenwichts vergelijking in termen van die alfa.
2:28:24
Zullen we nog eens doorgaan? Ik houd de alfa die daar staat.
2:28:34
Twee Mijn alfa gedeeld door een wortel, gedeeld door nog eens een wortel valt weg.
2:28:39
Er blijft één keer twee min alfa over.
2:28:44
Komt in de teller Als degene die daar staat dan heb ik een wortel van één mijn alfa en een één min als van tot de drie halven.
2:28:47
Dat is de één in alfa kwadraat die daar staat voor factor drie tot de drie halve
2:28:54
en dan heb ik een één op wortel twee en een één op twee tot de drie halve.
2:28:59
Dat is de factor vier die boven staat. Ik denk dat we er zijn. Wat doen we met de X?
2:29:04
Daarvoor moeten we terug naar thermodynamische tabellen.
2:29:13
En dan moet je altijd maar hopen dat de component waarin je geïnteresseerd bent dat die erin staat.
2:29:20
Maar ik verwacht dat ammoniak geen probleem gaat zijn.
2:29:25
Je ziet. Vormings vrije enthalpie van ammoniak -16.700.
2:29:28
Dus die X is een KP nul maal de exponentiële zorg is de KP nul is de exponentiële van
2:29:35
min delta één nul gedeeld door er twee maal p totaal op nul tot de macht min delta nu.
2:29:43
Als een wijnbar en kamertemperatuur werken, dan valt deze term weg.
2:29:55
Die gaan niet altijd weg vallen, want het beeld dat nu van de reactie is niet gelijk aan nul,
2:30:00
maar bij kamertemperatuur en één bar is wat er belangrijk is die exponentiële van min delta geen nul op RTE.
2:30:04
Minder dan geen nul op RTE 16.700.
2:30:14
Je ziet dat je krijgt exponentieel van zestien.
2:30:18
Duizendzevenhonderd gedeeld door RTE is ongeveer 3000.
2:30:21
Ga dus E tot de vijfde tot E tot de zesde zijn.
2:30:27
Dus je zit hier met een evenwicht constante die ongeveer gelijk is aan E tot de vijfde A tot de zesde.
2:30:31
Pak m beet ongeveer honderd gaat als zijn.
2:30:37
Dus je gaat om daar honderd te krijgen naar een limiet gaan waar die al voor in de richting van één moet gaan.
2:30:43
Je moet die teller naar nul krijgen om hoge evenwicht constante te hebben.
2:30:52
Dus bij kamertemperatuur en twee 98 Kelvin voorspellen we dat een evenwicht van die ammoniak vorming heel erg aan de ammoniak kant gaat liggen.
2:30:59
Je kan het keurig gaan uitwerken, want dit is nog altijd een kwadratische vergelijking,
2:31:08
maar je kan ook inschatten Om hier honderd te krijgen ga je een teller moeten hebben
2:31:12
die zal ga je noemen moeten hebben die naar nul gaat om het hele zaakje op te blazen.
2:31:16
Je alfa gaat naar één moeten gaan. Dus in principe mag er niks problematisch zijn met de reactie wat er niet gebeurd is.
2:31:21
Een feitelijke ammoniak vorming, dus daar zit wel een verschil tussen wat de thermodynamica voorspelt en wat er
2:31:37
op basis van de kinetiek op basis van de snelheid van een reactie mogelijk is.
2:31:43
Hier zit je met een reactie die eigenlijk een thermodynamisch gunstig is, maar omwille van kinetische redenen niet gaat gebeuren.
2:31:48
De figuur die je een hele tijd al linksboven gestaan heeft is Fritz Haber en dat is de man die hier ook bovenaan staat.
2:31:57
Haber Bosch proces. Wat Fritz Haber met die ammoniak synthese te maken heeft is dat hij omstandigheden
2:32:05
heeft ontwikkeld en een katalysator gevonden die de reactiesnelheid opdrijft.
2:32:11
Tenminste als je naar temperaturen gaat van ongeveer 400 à 500 graden Celsius.
2:32:19
En dus. Een katalysator is eigenlijk ijzerpoeder of ijzeroxide poeder en dus wat hij vaststelde is dat het
2:32:27
tempo waarmee ammoniak geproduceerd wordt omhoog gaat in de aanwezigheid van dat ijzerpoeder.
2:32:33
Dat is eigenlijk omdat stikstof en waterstof kunnen absorberen aan het oppervlak van dat ijzerpoeder.
2:32:39
En die dissociatie van stikstof moleculen is veel gemakkelijker in die geabsorbeerde toestand dan gewoon in de gasfase.
2:32:45
Dat is de rol van een katalysator. Dat geeft een snelle reactie pad om een omzetting te realiseren.
2:32:51
Als we gaan kijken naar de thermodynamica van die reactie. Ben daar de slides kwijt.
2:32:59
Als je kijkt naar de thermodynamica van die reactie, dan zie je dat de delta.
2:33:05
Dat die negatief is. Ik bedoel daarmee reactie.
2:33:11
Dat betekent dat we de temperatuur opdrijven en dan moet HBR doen omdat enkel bv 35 graden daar een zeker reactie tempo is.
2:33:17
Maar dan verschuift het evenwicht veeleer naar de kant van enkel stikstof en enkel waterstof.
2:33:26
Moest je die evenwichts constanten gaan berekenen bij 500 Kelvin,
2:33:32
dan zou je vinden dat dat geen tot de vijfde meer is, wat het ongeveer E tot een min vijfde gaat bedragen.
2:33:35
Dus in plaats van een evenwichts ligging helemaal rechts bij hogere temperaturen gaat het een heel erg lang worden die helemaal links komt te liggen.
2:33:40
Dat zie je ook in deze figuur. Als we werken bij een druk van één bar dan gaat Bij zelfs twintig graden Celsius gaat die evenwicht samenstelling.
2:33:47
In plaats van 50% ammoniak gaat hij teruggevallen zijn op maar een paar procenten ammoniak.
2:33:55
De reden waarom de naam van Bosch erbij staat is is ook een reactie met een beeld
2:34:03
dat nu gelijk aan -1 meer mooie stof aan de linkerkant dan aan de rechterkant.
2:34:08
Dat betekent dat je de druk opvoert. Dat je hevig kan gaan verschuiven van de kant van meer stof naar de kant van minder mooi stof.
2:34:14
En dus om die ammoniak synthese aan de praat te krijgen, hebben ze eigenlijk gebruik gemaakt van die inzichten in evenwichts thermodynamica.
2:34:25
We gaan naar hogere temperaturen om een zekere kinetiek te krijgen,
2:34:33
maar het verschuiven van de evenwichts ligging wordt teniet gedaan of wordt tegengewerkt door niet bij één bar te werken,
2:34:37
maar door bij drukken van 100 tot 1000 bar aan de slag te gaan.
2:34:44
En hier krijg je goede druk hoe de evenwichts samenstelling evolueert in functie van de druk.
2:34:48
En dan zie je die systematische toename van de fractie ammoniak.
2:34:53
Het is nog altijd niet 100% en je zit ergens aan 30% of 40% ammoniak bij stijgende druk.
2:34:57
Het is historisch gezien een belangrijke redactie om verschillende redenen.
2:35:04
En je hebt die oorspronkelijke motivatie om een soort ammoniak of stikstof industrie
2:35:11
mogelijk te maken zonder dat je die die guano moet gaan halen in de in de.
2:35:17
Ja, dat was de Stille Zuidzee. Dat is dus in die die eilanden die nogal ver van hier liggen.
2:35:23
Het resultaat van Haber en Bosh is gepubliceerd geworden.
2:35:29
Zijn patenten is gepatenteerd in duizendnegenhonderddertien. Is een belangrijke datum.
2:35:33
Duizendnegenhonderdveertien. Begint de Eerste Wereldoorlog.
2:35:43
Stikstof is niet alleen belangrijk voor kunstmest. Alle explosieven, alle ontploffingen, stoffen die gemaakt worden zijn stikstof componenten.
2:35:50
Omdat daar heel veel gasvormige stikstof bij geproduceerd wordt dat zeer sterk gebonden is.
2:36:00
Dus je gaat daar zeer veel energie bij vrijstellen. Tot de ontwikkeling van het Haber Bosch proces Was die hele dynamiet en buskruit en en, en en en.
2:36:06
Explosieve industrie was ook gebaseerd op die guano die uit die eilanden van de Stille Zuidzee moest komen.
2:36:21
Als je naar kaarten gaat kijken van de oorlogs activiteiten in het begin van de eerste wereldoorlog,
2:36:28
dan ga je zien dat er rond die eilanden dat daar ook slag geleverd is.
2:36:36
Daar zijn zeeslagen rond geleverd om die controle over die eilanden te krijgen en de de Duitse zeemacht heeft
2:36:40
die verloren en dus van in het begin van de Eerste Wereldoorlog had Duitsland geen toegang tot die nitraten.
2:36:48
Konden ze dus ook geen. Buskruit als het ware maken.
2:36:56
Ware het niet dat in duizendnegenhonderddertien dat Haber Bosch proces ontwikkeld was,
2:37:03
dat eigenlijk de Duitse industrie vier jaar lang in staat gesteld heeft om die oorlogsinspanning gaande te houden.
2:37:07
Er is dus een chemie die opgebouwd is vanuit het idee we moeten de mensheid van van voedsel voorzien.
2:37:14
Is eigenlijk in de praktijk gebruik geworden om vier jaar lang oorlog te voeren.
2:37:20
Die Fritz Haber is is in een een, een, een. Veel, kan het veelzijdig of is een figuur met zeer veel facetten.
2:37:25
Heeft ook in duizendnegenhonderdvijftien een brief aan de Duitse keizer gestuurd Fritz Haber,
2:37:32
waarin hij uitlegt dat volgens hem chloorgas een zeer interessant middel is in de in de oorlogsvoering.
2:37:37
Het is de man die eigenlijk al die gifgas aanvallen. Het initiatief daarvoor genomen heeft.
2:37:46
Bij de eerste gifgas was die ook aanwezig om meer te zeggen van kijk die chloorgas potten,
2:37:53
die moeten daar en daar en daar opgesteld staan opdat dat chloor niet in de Duitse linies, maar in de in de in de Engelse linies zou terechtkomen.
2:37:58
Heeft in 1919 Op het einde of na de Tweede Eerste Wereldoorlog heeft hij de Nobelprijs Chemie gekregen,
2:38:07
exact voor de ontwikkeling van die ammoniak synthese dus, is een figuur waar heel veel zaken bij samenkomen.
2:38:13
Hij was ook Joods en heeft in. Hij was dan aangesteld als hoofd van het Fritz Haber Instituut in Berlijn en dan,
2:38:21
in duizendnegenhonderddrieëndertig, moest hij al zijn werknemers ontslaan.
2:38:28
Op dat moment was de nazipartij aan de macht gekomen.
2:38:32
En die, die vonden het niet goed dat er Joodse wetenschappers waren en hij werd dus verplicht om iedereen te ontslaan, eigenlijk zichzelf inbegrepen,
2:38:35
heeft het dan geweigerd, heeft dat niet willen doen en uiteindelijk heeft hij dan geprobeerd asiel te krijgen in het Verenigd Koninkrijk,
2:38:43
maar die kenden hem natuurlijk van ja, de chemicus die de Eerste Wereldoorlog vier jaar lang heeft laten functioneren en ze hebben
2:38:50
hem dat asiel geweigerd hebben dat niet gegeven is dan uiteindelijk in 1934 overleden.
2:38:57
Is dus een hele complexe figuur. Je moet daar iets over lezen omdat daar eigenlijk, ja,
2:39:03
zowel het goede als het slechte van de wetenschap in samenkomt en en uitlegt dat dat de wereld vaak ingewikkeld in elkaar zit.
2:39:07
Goed, dat is wat ik er wil over zeggen, maar is een hele mooie toepassing,
2:39:16
ook voor ons doeleinde dan van die verschuivingen van chemische evenwichten onder invloed van temperatuur en druk.
2:39:19

\end{document}